\chapter{睡眠与唤醒}
\label{CH:SLEEP}
%% 

调度和锁有助于将一个线程的动作对其他线程隐藏起来，
但我们也需要一些抽象来帮助线程之间有意图地进行交互。
例如，xv6 中管道的读者可能需要等待写进程产生数据；
父进程对 \lstinline{wait} 的调用可能需要等待子进程退出；
而读取磁盘的进程则需要等待磁盘硬件完成读取操作。
在这些情况（以及许多其他情况）下，xv6 内核使用一种称为睡眠和唤醒的机制。
睡眠允许一个内核线程等待某个条件变为真；
另一个线程或一个中断处理程序可以使该条件变为真（通常通过修改某些变量），
然后调用唤醒，以指示等待该条件的线程应恢复执行。
睡眠和唤醒通常被称为
\indextext{顺序协调}
或
\indextext{条件同步}
机制。

在继续阅读之前，请先阅读 {\tt kernel/proc.c} 中的函数 {\tt sleep()} 和 {\tt
  wakeup()}，以及整个文件 {\tt kernel/pipe.c}。

\section{概述}

睡眠/唤醒接口如下所示：

\begin{lstlisting}
void sleep(void *chan, struct spinlock *lk)
void wakeup(void *chan)
\end{lstlisting}

\lstinline{sleep()} 将调用进程标记为 \lstinline{SLEEPING}（而不是 \lstinline{RUNNABLE}），
并通过上下文切换到调度器来释放 CPU，以便其他进程可以运行。
\lstinline{chan} 参数称为\indextext{等待通道}。
\lstinline{wakeup(chan)} 会唤醒所有（如果有的话）使用相同 \lstinline{chan}
值调用过 \lstinline{sleep(chan, ...)} 的进程。
\lstinline{sleep} 和 \lstinline{wakeup} 将 \lstinline{chan}
视为一个不透明的 64 位值；它们对该值所做的唯一操作就是进行相等性比较。
通常的模式是调用者传递某个方便对象的地址作为 \lstinline{chan} 参数。

内核代码调用 \lstinline{sleep} 来等待某个
\indextext{条件}
变为真。
例如，从管道读取数据的内核代码在管道缓冲区当前为空时会调用 \lstinline{sleep}；
这种情况下的条件是管道缓冲区变为非空（由于另一个进程向管道写入数据）。
\lstinline{sleep} 和 \lstinline{wakeup} 并不知道条件是什么：
只有调用代码知道。
通常的模式是调用者首先检查条件，如果条件不成立则调用 \lstinline{sleep}；
之后使条件变为真的代码会调用 \lstinline{wakeup}。

下面简要展示了 xv6 内核管道代码如何使用
\lstinline{sleep} 和 \lstinline{wakeup}：

\begin{lstlisting}
piperead(pipe){
  acquire(&pipe->lock);
  while(there's no data in pipe->buffer){
    (*@\textcolor{blue}{// ZZZ}@*)
    sleep(&pipe, &pipe->lock);
  }
  remove the data from the pipe;
  release(&pipe->lock);
}

pipewrite(pipe){
  acquire(&pipe->lock);
  append data to pipe->buffer;
  wakeup(&pipe);
  release(&pipe->lock);
}
\end{lstlisting}

这段代码使用管道数据结构的地址作为等待通道。

\lstinline{sleep} 的 \lstinline{lk} 参数是什么？
在所有使用 \lstinline{sleep}/\lstinline{wakeup} 的场景中，
条件都涉及共享数据，由睡眠线程和调用 \lstinline{wakeup} 的线程共同使用，
因此总是存在一个锁来保护该条件。这个锁被称为\indextext{条件锁}。
在上面的管道代码中，两个函数都在持有管道锁的情况下使用管道及其缓冲区，
在这种情况下，管道锁也是条件锁。
规则是，任何调用 \lstinline{sleep} 或 \lstinline{wakeup} 的代码都必须持有条件锁，
并且该锁必须作为第二个参数传递给 \lstinline{sleep}。

在调用 \lstinline{sleep} 时必须持有条件锁，并将其传递给 \lstinline{sleep}，
其原因是为了防止另一个线程可能在检查条件之后、调用 \lstinline{sleep} 之前
调用 \lstinline{wakeup}。
此时调用 \lstinline{wakeup} 会发现没有正在睡眠的进程可以唤醒；
\lstinline{wakeup} 就会直接返回。
但随后对 \lstinline{sleep} 的调用可能永远不会被唤醒，
因为本应唤醒它的 \lstinline{wakeup} 已经发生过了。
这种不良情况被称为\indextext{丢失唤醒}。

在上面的管道示例中，所要避免的丢失唤醒是：
另一个 CPU 上的线程可能在 {\tt ZZZ} 标记点、
即在 \lstinline{piperead} 检查条件之后、调用 \lstinline{sleep} 之前，
调用 \lstinline{pipewrite}。
\lstinline{piperead} 在检查条件到调用 \lstinline{sleep} 这段时间内持有管道锁，
这一事实阻止了 \lstinline{pipewrite} 的执行，从而防止了丢失唤醒。

{\tt sleep()} 会释放条件锁，以便调用 {\tt wakeup()} 的代码可以继续执行。
{\tt sleep()} 还会上下文切换到调度器，以便在等待期间让其他线程运行。
其实现以相对于 {\tt wakeup()} 原子（不可分割）的方式执行这两个步骤，
以防止丢失唤醒。

%% 
\section{代码：睡眠与唤醒}
%% 

Xv6 的
\indexcode{sleep}
\lineref{kernel/proc.c:/^sleep/}
和
\indexcode{wakeup}
\lineref{kernel/proc.c:/^wakeup/}
实现了上述示例中使用的接口。
基本思想是让
\lstinline{sleep}
将当前进程标记为
\indexcode{SLEEPING}，
然后调用
\indexcode{sched}
以释放 CPU；
\lstinline{wakeup}
则查找在给定等待通道上睡眠的进程，
并将其标记为
\indexcode{RUNNABLE}。

\lstinline{sleep}
首先获取
\indexcode{p->lock}
\lineref{kernel/proc.c:/DOC: sleeplock1/}，
然后{\it 才}释放条件锁
\lstinline{lk}。
\lstinline{sleep}
在任何时候都持有这两个锁中的一个或两个，
这一事实阻止了并发的 \lstinline{wakeup}（它必须获取并同时持有这两个锁）采取行动，
从而防止了丢失唤醒。
现在，
\lstinline{sleep}
只持有
\lstinline{p->lock}，
它可以通过记录等待通道、
将进程状态更改为 \texttt{SLEEPING}、
并调用
\lstinline{sched}
\linerefs{kernel/proc.c:/chan.=.chan/,/sched/}
来使进程进入睡眠状态。
稍后会明白，为什么在进程被标记为 \texttt{SLEEPING} 之前不释放 \lstinline{p->lock}（由调度器释放）至关重要。

在某个时刻，一个进程将获取条件锁，
设置睡眠者正在等待的条件，
并调用 \lstinline{wakeup(chan)}。
重要的是，\lstinline{wakeup} 必须在持有条件锁的情况下被调用\footnote{%
%
严格来说，只要 \lstinline{wakeup} 在 \lstinline{acquire} 之后即可
（也就是说，可以在 \lstinline{release} 之后调用 \lstinline{wakeup}）。%
%
}。
\lstinline{wakeup}
遍历进程表
\lineref{kernel/proc.c:/^wakeup\(/}。
它会获取所检查的每个进程的 \lstinline{p->lock}。
当 \lstinline{wakeup} 找到一个状态为
\indexcode{SLEEPING} 且
\indexcode{chan} 匹配的进程时，
它将把该进程的状态更改为
\indexcode{RUNNABLE}。
下次 \lstinline{scheduler} 运行时，
它将看到该进程已准备好运行。

\begin{figure}[t]
  \chapter{睡眠与唤醒}
\label{CH:SLEEP}
%% 

调度和锁有助于将一个线程的动作对其他线程隐藏起来，
但我们也需要一些抽象来帮助线程之间有意图地进行交互。
例如，xv6 中管道的读者可能需要等待写进程产生数据；
父进程对 \lstinline{wait} 的调用可能需要等待子进程退出；
而读取磁盘的进程则需要等待磁盘硬件完成读取操作。
在这些情况（以及许多其他情况）下，xv6 内核使用一种称为睡眠和唤醒的机制。
睡眠允许一个内核线程等待某个条件变为真；
另一个线程或一个中断处理程序可以使该条件变为真（通常通过修改某些变量），
然后调用唤醒，以指示等待该条件的线程应恢复执行。
睡眠和唤醒通常被称为
\indextext{顺序协调}
或
\indextext{条件同步}
机制。

在继续阅读之前，请先阅读 {\tt kernel/proc.c} 中的函数 {\tt sleep()} 和 {\tt
  wakeup()}，以及整个文件 {\tt kernel/pipe.c}。

\section{概述}

睡眠/唤醒接口如下所示：

\begin{lstlisting}
void sleep(void *chan, struct spinlock *lk)
void wakeup(void *chan)
\end{lstlisting}

\lstinline{sleep()} 将调用进程标记为 \lstinline{SLEEPING}（而不是 \lstinline{RUNNABLE}），
并通过上下文切换到调度器来释放 CPU，以便其他进程可以运行。
\lstinline{chan} 参数称为\indextext{等待通道}。
\lstinline{wakeup(chan)} 会唤醒所有（如果有的话）使用相同 \lstinline{chan}
值调用过 \lstinline{sleep(chan, ...)} 的进程。
\lstinline{sleep} 和 \lstinline{wakeup} 将 \lstinline{chan}
视为一个不透明的 64 位值；它们对该值所做的唯一操作就是进行相等性比较。
通常的模式是调用者传递某个方便对象的地址作为 \lstinline{chan} 参数。

内核代码调用 \lstinline{sleep} 来等待某个
\indextext{条件}
变为真。
例如，从管道读取数据的内核代码在管道缓冲区当前为空时会调用 \lstinline{sleep}；
这种情况下的条件是管道缓冲区变为非空（由于另一个进程向管道写入数据）。
\lstinline{sleep} 和 \lstinline{wakeup} 并不知道条件是什么：
只有调用代码知道。
通常的模式是调用者首先检查条件，如果条件不成立则调用 \lstinline{sleep}；
之后使条件变为真的代码会调用 \lstinline{wakeup}。

下面简要展示了 xv6 内核管道代码如何使用
\lstinline{sleep} 和 \lstinline{wakeup}：

\begin{lstlisting}
piperead(pipe){
  acquire(&pipe->lock);
  while(there's no data in pipe->buffer){
    (*@\textcolor{blue}{// ZZZ}@*)
    sleep(&pipe, &pipe->lock);
  }
  remove the data from the pipe;
  release(&pipe->lock);
}

pipewrite(pipe){
  acquire(&pipe->lock);
  append data to pipe->buffer;
  wakeup(&pipe);
  release(&pipe->lock);
}
\end{lstlisting}

这段代码使用管道数据结构的地址作为等待通道。

\lstinline{sleep} 的 \lstinline{lk} 参数是什么？
在所有使用 \lstinline{sleep}/\lstinline{wakeup} 的场景中，
条件都涉及共享数据，由睡眠线程和调用 \lstinline{wakeup} 的线程共同使用，
因此总是存在一个锁来保护该条件。这个锁被称为\indextext{条件锁}。
在上面的管道代码中，两个函数都在持有管道锁的情况下使用管道及其缓冲区，
在这种情况下，管道锁也是条件锁。
规则是，任何调用 \lstinline{sleep} 或 \lstinline{wakeup} 的代码都必须持有条件锁，
并且该锁必须作为第二个参数传递给 \lstinline{sleep}。

在调用 \lstinline{sleep} 时必须持有条件锁，并将其传递给 \lstinline{sleep}，
其原因是为了防止另一个线程可能在检查条件之后、调用 \lstinline{sleep} 之前
调用 \lstinline{wakeup}。
此时调用 \lstinline{wakeup} 会发现没有正在睡眠的进程可以唤醒；
\lstinline{wakeup} 就会直接返回。
但随后对 \lstinline{sleep} 的调用可能永远不会被唤醒，
因为本应唤醒它的 \lstinline{wakeup} 已经发生过了。
这种不良情况被称为\indextext{丢失唤醒}。

在上面的管道示例中，所要避免的丢失唤醒是：
另一个 CPU 上的线程可能在 {\tt ZZZ} 标记点、
即在 \lstinline{piperead} 检查条件之后、调用 \lstinline{sleep} 之前，
调用 \lstinline{pipewrite}。
\lstinline{piperead} 在检查条件到调用 \lstinline{sleep} 这段时间内持有管道锁，
这一事实阻止了 \lstinline{pipewrite} 的执行，从而防止了丢失唤醒。

{\tt sleep()} 会释放条件锁，以便调用 {\tt wakeup()} 的代码可以继续执行。
{\tt sleep()} 还会上下文切换到调度器，以便在等待期间让其他线程运行。
其实现以相对于 {\tt wakeup()} 原子（不可分割）的方式执行这两个步骤，
以防止丢失唤醒。

%% 
\section{代码：睡眠与唤醒}
%% 

Xv6 的
\indexcode{sleep}
\lineref{kernel/proc.c:/^sleep/}
和
\indexcode{wakeup}
\lineref{kernel/proc.c:/^wakeup/}
实现了上述示例中使用的接口。
基本思想是让
\lstinline{sleep}
将当前进程标记为
\indexcode{SLEEPING}，
然后调用
\indexcode{sched}
以释放 CPU；
\lstinline{wakeup}
则查找在给定等待通道上睡眠的进程，
并将其标记为
\indexcode{RUNNABLE}。

\lstinline{sleep}
首先获取
\indexcode{p->lock}
\lineref{kernel/proc.c:/DOC: sleeplock1/}，
然后{\it 才}释放条件锁
\lstinline{lk}。
\lstinline{sleep}
在任何时候都持有这两个锁中的一个或两个，
这一事实阻止了并发的 \lstinline{wakeup}（它必须获取并同时持有这两个锁）采取行动，
从而防止了丢失唤醒。
现在，
\lstinline{sleep}
只持有
\lstinline{p->lock}，
它可以通过记录等待通道、
将进程状态更改为 \texttt{SLEEPING}、
并调用
\lstinline{sched}
\linerefs{kernel/proc.c:/chan.=.chan/,/sched/}
来使进程进入睡眠状态。
稍后会明白，为什么在进程被标记为 \texttt{SLEEPING} 之前不释放 \lstinline{p->lock}（由调度器释放）至关重要。

在某个时刻，一个进程将获取条件锁，
设置睡眠者正在等待的条件，
并调用 \lstinline{wakeup(chan)}。
重要的是，\lstinline{wakeup} 必须在持有条件锁的情况下被调用\footnote{%
%
严格来说，只要 \lstinline{wakeup} 在 \lstinline{acquire} 之后即可
（也就是说，可以在 \lstinline{release} 之后调用 \lstinline{wakeup}）。%
%
}。
\lstinline{wakeup}
遍历进程表
\lineref{kernel/proc.c:/^wakeup\(/}。
它会获取所检查的每个进程的 \lstinline{p->lock}。
当 \lstinline{wakeup} 找到一个状态为
\indexcode{SLEEPING} 且
\indexcode{chan} 匹配的进程时，
它将把该进程的状态更改为
\indexcode{RUNNABLE}。
下次 \lstinline{scheduler} 运行时，
它将看到该进程已准备好运行。

\begin{figure}[t]
  \chapter{睡眠与唤醒}
\label{CH:SLEEP}
%% 

调度和锁有助于将一个线程的动作对其他线程隐藏起来，
但我们也需要一些抽象来帮助线程之间有意图地进行交互。
例如，xv6 中管道的读者可能需要等待写进程产生数据；
父进程对 \lstinline{wait} 的调用可能需要等待子进程退出；
而读取磁盘的进程则需要等待磁盘硬件完成读取操作。
在这些情况（以及许多其他情况）下，xv6 内核使用一种称为睡眠和唤醒的机制。
睡眠允许一个内核线程等待某个条件变为真；
另一个线程或一个中断处理程序可以使该条件变为真（通常通过修改某些变量），
然后调用唤醒，以指示等待该条件的线程应恢复执行。
睡眠和唤醒通常被称为
\indextext{顺序协调}
或
\indextext{条件同步}
机制。

在继续阅读之前，请先阅读 {\tt kernel/proc.c} 中的函数 {\tt sleep()} 和 {\tt
  wakeup()}，以及整个文件 {\tt kernel/pipe.c}。

\section{概述}

睡眠/唤醒接口如下所示：

\begin{lstlisting}
void sleep(void *chan, struct spinlock *lk)
void wakeup(void *chan)
\end{lstlisting}

\lstinline{sleep()} 将调用进程标记为 \lstinline{SLEEPING}（而不是 \lstinline{RUNNABLE}），
并通过上下文切换到调度器来释放 CPU，以便其他进程可以运行。
\lstinline{chan} 参数称为\indextext{等待通道}。
\lstinline{wakeup(chan)} 会唤醒所有（如果有的话）使用相同 \lstinline{chan}
值调用过 \lstinline{sleep(chan, ...)} 的进程。
\lstinline{sleep} 和 \lstinline{wakeup} 将 \lstinline{chan}
视为一个不透明的 64 位值；它们对该值所做的唯一操作就是进行相等性比较。
通常的模式是调用者传递某个方便对象的地址作为 \lstinline{chan} 参数。

内核代码调用 \lstinline{sleep} 来等待某个
\indextext{条件}
变为真。
例如，从管道读取数据的内核代码在管道缓冲区当前为空时会调用 \lstinline{sleep}；
这种情况下的条件是管道缓冲区变为非空（由于另一个进程向管道写入数据）。
\lstinline{sleep} 和 \lstinline{wakeup} 并不知道条件是什么：
只有调用代码知道。
通常的模式是调用者首先检查条件，如果条件不成立则调用 \lstinline{sleep}；
之后使条件变为真的代码会调用 \lstinline{wakeup}。

下面简要展示了 xv6 内核管道代码如何使用
\lstinline{sleep} 和 \lstinline{wakeup}：

\begin{lstlisting}
piperead(pipe){
  acquire(&pipe->lock);
  while(there's no data in pipe->buffer){
    (*@\textcolor{blue}{// ZZZ}@*)
    sleep(&pipe, &pipe->lock);
  }
  remove the data from the pipe;
  release(&pipe->lock);
}

pipewrite(pipe){
  acquire(&pipe->lock);
  append data to pipe->buffer;
  wakeup(&pipe);
  release(&pipe->lock);
}
\end{lstlisting}

这段代码使用管道数据结构的地址作为等待通道。

\lstinline{sleep} 的 \lstinline{lk} 参数是什么？
在所有使用 \lstinline{sleep}/\lstinline{wakeup} 的场景中，
条件都涉及共享数据，由睡眠线程和调用 \lstinline{wakeup} 的线程共同使用，
因此总是存在一个锁来保护该条件。这个锁被称为\indextext{条件锁}。
在上面的管道代码中，两个函数都在持有管道锁的情况下使用管道及其缓冲区，
在这种情况下，管道锁也是条件锁。
规则是，任何调用 \lstinline{sleep} 或 \lstinline{wakeup} 的代码都必须持有条件锁，
并且该锁必须作为第二个参数传递给 \lstinline{sleep}。

在调用 \lstinline{sleep} 时必须持有条件锁，并将其传递给 \lstinline{sleep}，
其原因是为了防止另一个线程可能在检查条件之后、调用 \lstinline{sleep} 之前
调用 \lstinline{wakeup}。
此时调用 \lstinline{wakeup} 会发现没有正在睡眠的进程可以唤醒；
\lstinline{wakeup} 就会直接返回。
但随后对 \lstinline{sleep} 的调用可能永远不会被唤醒，
因为本应唤醒它的 \lstinline{wakeup} 已经发生过了。
这种不良情况被称为\indextext{丢失唤醒}。

在上面的管道示例中，所要避免的丢失唤醒是：
另一个 CPU 上的线程可能在 {\tt ZZZ} 标记点、
即在 \lstinline{piperead} 检查条件之后、调用 \lstinline{sleep} 之前，
调用 \lstinline{pipewrite}。
\lstinline{piperead} 在检查条件到调用 \lstinline{sleep} 这段时间内持有管道锁，
这一事实阻止了 \lstinline{pipewrite} 的执行，从而防止了丢失唤醒。

{\tt sleep()} 会释放条件锁，以便调用 {\tt wakeup()} 的代码可以继续执行。
{\tt sleep()} 还会上下文切换到调度器，以便在等待期间让其他线程运行。
其实现以相对于 {\tt wakeup()} 原子（不可分割）的方式执行这两个步骤，
以防止丢失唤醒。

%% 
\section{代码：睡眠与唤醒}
%% 

Xv6 的
\indexcode{sleep}
\lineref{kernel/proc.c:/^sleep/}
和
\indexcode{wakeup}
\lineref{kernel/proc.c:/^wakeup/}
实现了上述示例中使用的接口。
基本思想是让
\lstinline{sleep}
将当前进程标记为
\indexcode{SLEEPING}，
然后调用
\indexcode{sched}
以释放 CPU；
\lstinline{wakeup}
则查找在给定等待通道上睡眠的进程，
并将其标记为
\indexcode{RUNNABLE}。

\lstinline{sleep}
首先获取
\indexcode{p->lock}
\lineref{kernel/proc.c:/DOC: sleeplock1/}，
然后{\it 才}释放条件锁
\lstinline{lk}。
\lstinline{sleep}
在任何时候都持有这两个锁中的一个或两个，
这一事实阻止了并发的 \lstinline{wakeup}（它必须获取并同时持有这两个锁）采取行动，
从而防止了丢失唤醒。
现在，
\lstinline{sleep}
只持有
\lstinline{p->lock}，
它可以通过记录等待通道、
将进程状态更改为 \texttt{SLEEPING}、
并调用
\lstinline{sched}
\linerefs{kernel/proc.c:/chan.=.chan/,/sched/}
来使进程进入睡眠状态。
稍后会明白，为什么在进程被标记为 \texttt{SLEEPING} 之前不释放 \lstinline{p->lock}（由调度器释放）至关重要。

在某个时刻，一个进程将获取条件锁，
设置睡眠者正在等待的条件，
并调用 \lstinline{wakeup(chan)}。
重要的是，\lstinline{wakeup} 必须在持有条件锁的情况下被调用\footnote{%
%
严格来说，只要 \lstinline{wakeup} 在 \lstinline{acquire} 之后即可
（也就是说，可以在 \lstinline{release} 之后调用 \lstinline{wakeup}）。%
%
}。
\lstinline{wakeup}
遍历进程表
\lineref{kernel/proc.c:/^wakeup\(/}。
它会获取所检查的每个进程的 \lstinline{p->lock}。
当 \lstinline{wakeup} 找到一个状态为
\indexcode{SLEEPING} 且
\indexcode{chan} 匹配的进程时，
它将把该进程的状态更改为
\indexcode{RUNNABLE}。
下次 \lstinline{scheduler} 运行时，
它将看到该进程已准备好运行。

\begin{figure}[t]
  \chapter{睡眠与唤醒}
\label{CH:SLEEP}
%% 

调度和锁有助于将一个线程的动作对其他线程隐藏起来，
但我们也需要一些抽象来帮助线程之间有意图地进行交互。
例如，xv6 中管道的读者可能需要等待写进程产生数据；
父进程对 \lstinline{wait} 的调用可能需要等待子进程退出；
而读取磁盘的进程则需要等待磁盘硬件完成读取操作。
在这些情况（以及许多其他情况）下，xv6 内核使用一种称为睡眠和唤醒的机制。
睡眠允许一个内核线程等待某个条件变为真；
另一个线程或一个中断处理程序可以使该条件变为真（通常通过修改某些变量），
然后调用唤醒，以指示等待该条件的线程应恢复执行。
睡眠和唤醒通常被称为
\indextext{顺序协调}
或
\indextext{条件同步}
机制。

在继续阅读之前，请先阅读 {\tt kernel/proc.c} 中的函数 {\tt sleep()} 和 {\tt
  wakeup()}，以及整个文件 {\tt kernel/pipe.c}。

\section{概述}

睡眠/唤醒接口如下所示：

\begin{lstlisting}
void sleep(void *chan, struct spinlock *lk)
void wakeup(void *chan)
\end{lstlisting}

\lstinline{sleep()} 将调用进程标记为 \lstinline{SLEEPING}（而不是 \lstinline{RUNNABLE}），
并通过上下文切换到调度器来释放 CPU，以便其他进程可以运行。
\lstinline{chan} 参数称为\indextext{等待通道}。
\lstinline{wakeup(chan)} 会唤醒所有（如果有的话）使用相同 \lstinline{chan}
值调用过 \lstinline{sleep(chan, ...)} 的进程。
\lstinline{sleep} 和 \lstinline{wakeup} 将 \lstinline{chan}
视为一个不透明的 64 位值；它们对该值所做的唯一操作就是进行相等性比较。
通常的模式是调用者传递某个方便对象的地址作为 \lstinline{chan} 参数。

内核代码调用 \lstinline{sleep} 来等待某个
\indextext{条件}
变为真。
例如，从管道读取数据的内核代码在管道缓冲区当前为空时会调用 \lstinline{sleep}；
这种情况下的条件是管道缓冲区变为非空（由于另一个进程向管道写入数据）。
\lstinline{sleep} 和 \lstinline{wakeup} 并不知道条件是什么：
只有调用代码知道。
通常的模式是调用者首先检查条件，如果条件不成立则调用 \lstinline{sleep}；
之后使条件变为真的代码会调用 \lstinline{wakeup}。

下面简要展示了 xv6 内核管道代码如何使用
\lstinline{sleep} 和 \lstinline{wakeup}：

\begin{lstlisting}
piperead(pipe){
  acquire(&pipe->lock);
  while(there's no data in pipe->buffer){
    (*@\textcolor{blue}{// ZZZ}@*)
    sleep(&pipe, &pipe->lock);
  }
  remove the data from the pipe;
  release(&pipe->lock);
}

pipewrite(pipe){
  acquire(&pipe->lock);
  append data to pipe->buffer;
  wakeup(&pipe);
  release(&pipe->lock);
}
\end{lstlisting}

这段代码使用管道数据结构的地址作为等待通道。

\lstinline{sleep} 的 \lstinline{lk} 参数是什么？
在所有使用 \lstinline{sleep}/\lstinline{wakeup} 的场景中，
条件都涉及共享数据，由睡眠线程和调用 \lstinline{wakeup} 的线程共同使用，
因此总是存在一个锁来保护该条件。这个锁被称为\indextext{条件锁}。
在上面的管道代码中，两个函数都在持有管道锁的情况下使用管道及其缓冲区，
在这种情况下，管道锁也是条件锁。
规则是，任何调用 \lstinline{sleep} 或 \lstinline{wakeup} 的代码都必须持有条件锁，
并且该锁必须作为第二个参数传递给 \lstinline{sleep}。

在调用 \lstinline{sleep} 时必须持有条件锁，并将其传递给 \lstinline{sleep}，
其原因是为了防止另一个线程可能在检查条件之后、调用 \lstinline{sleep} 之前
调用 \lstinline{wakeup}。
此时调用 \lstinline{wakeup} 会发现没有正在睡眠的进程可以唤醒；
\lstinline{wakeup} 就会直接返回。
但随后对 \lstinline{sleep} 的调用可能永远不会被唤醒，
因为本应唤醒它的 \lstinline{wakeup} 已经发生过了。
这种不良情况被称为\indextext{丢失唤醒}。

在上面的管道示例中，所要避免的丢失唤醒是：
另一个 CPU 上的线程可能在 {\tt ZZZ} 标记点、
即在 \lstinline{piperead} 检查条件之后、调用 \lstinline{sleep} 之前，
调用 \lstinline{pipewrite}。
\lstinline{piperead} 在检查条件到调用 \lstinline{sleep} 这段时间内持有管道锁，
这一事实阻止了 \lstinline{pipewrite} 的执行，从而防止了丢失唤醒。

{\tt sleep()} 会释放条件锁，以便调用 {\tt wakeup()} 的代码可以继续执行。
{\tt sleep()} 还会上下文切换到调度器，以便在等待期间让其他线程运行。
其实现以相对于 {\tt wakeup()} 原子（不可分割）的方式执行这两个步骤，
以防止丢失唤醒。

%% 
\section{代码：睡眠与唤醒}
%% 

Xv6 的
\indexcode{sleep}
\lineref{kernel/proc.c:/^sleep/}
和
\indexcode{wakeup}
\lineref{kernel/proc.c:/^wakeup/}
实现了上述示例中使用的接口。
基本思想是让
\lstinline{sleep}
将当前进程标记为
\indexcode{SLEEPING}，
然后调用
\indexcode{sched}
以释放 CPU；
\lstinline{wakeup}
则查找在给定等待通道上睡眠的进程，
并将其标记为
\indexcode{RUNNABLE}。

\lstinline{sleep}
首先获取
\indexcode{p->lock}
\lineref{kernel/proc.c:/DOC: sleeplock1/}，
然后{\it 才}释放条件锁
\lstinline{lk}。
\lstinline{sleep}
在任何时候都持有这两个锁中的一个或两个，
这一事实阻止了并发的 \lstinline{wakeup}（它必须获取并同时持有这两个锁）采取行动，
从而防止了丢失唤醒。
现在，
\lstinline{sleep}
只持有
\lstinline{p->lock}，
它可以通过记录等待通道、
将进程状态更改为 \texttt{SLEEPING}、
并调用
\lstinline{sched}
\linerefs{kernel/proc.c:/chan.=.chan/,/sched/}
来使进程进入睡眠状态。
稍后会明白，为什么在进程被标记为 \texttt{SLEEPING} 之前不释放 \lstinline{p->lock}（由调度器释放）至关重要。

在某个时刻，一个进程将获取条件锁，
设置睡眠者正在等待的条件，
并调用 \lstinline{wakeup(chan)}。
重要的是，\lstinline{wakeup} 必须在持有条件锁的情况下被调用\footnote{%
%
严格来说，只要 \lstinline{wakeup} 在 \lstinline{acquire} 之后即可
（也就是说，可以在 \lstinline{release} 之后调用 \lstinline{wakeup}）。%
%
}。
\lstinline{wakeup}
遍历进程表
\lineref{kernel/proc.c:/^wakeup\(/}。
它会获取所检查的每个进程的 \lstinline{p->lock}。
当 \lstinline{wakeup} 找到一个状态为
\indexcode{SLEEPING} 且
\indexcode{chan} 匹配的进程时，
它将把该进程的状态更改为
\indexcode{RUNNABLE}。
下次 \lstinline{scheduler} 运行时，
它将看到该进程已准备好运行。

\begin{figure}[t]
  \input{fig/sleep.tex}
  \caption{重叠锁以避免丢失唤醒}
    \label{fig:overlap}
\end{figure}

为什么 \lstinline{sleep} 和 \lstinline{wakeup} 的锁规则能确保即将进入睡眠的进程不会错过并发的唤醒？
即将进入睡眠的进程从{\it 在}检查条件之前，到{\it 在}将其自身标记为 \texttt{SLEEPING} 之后，
始终持有条件锁或其自身的 \lstinline{p->lock} 或两者都持有；
见图~\ref{fig:overlap}。
调用 \texttt{wakeup} 的进程需要获取{\it 两个}锁。
唤醒者可能先获取锁，这意味着它会在消费者线程检查条件之前使条件成立，
这样消费者线程就不需要调用 {\tt sleep()}；
或者唤醒者的 {\tt acquire()} 可能需要等待，直到消费者线程完全完成进入睡眠并释放锁，
此时唤醒者将看到消费者线程已被标记为 {\tt SLEEPING} 并将其唤醒。

有时多个进程在同一个通道上睡眠；例如，多个进程从一个管道读取数据。
一次对 \lstinline{wakeup} 的调用将唤醒它们全部。
其中一个将首先运行，并获取 \lstinline{sleep} 调用时所持有的锁，
并且（在管道的例子中）读取所有等待的数据。
其他进程会发现，尽管被唤醒了，但没有数据可读。
从它们的角度来看，这次唤醒是“虚假的”，它们必须再次进入睡眠。
因此，\lstinline{sleep} 总是在一个重新检查条件的循环内被调用，如上面的 \lstinline{P} 所示。

如果两个使用 sleep/wakeup 的地方意外选择了相同的通道，也不会造成伤害：
它们会遇到虚假唤醒，但如上所述的循环可以容忍这个问题。
sleep/wakeup 的魅力很大程度上在于它既轻量级（无需创建特殊的数据结构作为等待通道），
又提供了一层间接性（调用者无需知道他们正在与哪个特定进程交互）。

%% 
\section{代码：管道}
%% 
xv6 的管道是使用 \lstinline{sleep} 和 \lstinline{wakeup} 来同步生产者和消费者的代码示例。
我们在第~\ref{CH:UNIX}~章中了解了管道的接口：
写入管道一端的数据被复制到内核缓冲区，然后可以从另一端读取。
后续章节将研究围绕管道的文件描述符支持，现在让我们看看
\indexcode{pipewrite}
和
\indexcode{piperead}
的实现。

每个管道由一个
\indexcode{struct pipe}
表示，其中包含一个
\lstinline{lock}
和一个
\lstinline{data}
缓冲区。
字段
\lstinline{nread}
和
\lstinline{nwrite}
统计从缓冲区读取和写入缓冲区的总字节数。
缓冲区是环绕的：
在
\lstinline{buf[PIPESIZE-1]}
之后写入的下一个字节是
\lstinline{buf[0]}。
计数不会环绕。
这个约定使得实现能够区分缓冲区满
（\lstinline{nwrite}
\lstinline{==}
\lstinline{nread+PIPESIZE}）
和缓冲区空
（\lstinline{nwrite}
\lstinline{==}
\lstinline{nread}），
但这意味着对缓冲区的索引必须使用
\lstinline{buf[nread}
\lstinline{%}
\lstinline{PIPESIZE]}
而不是简单的
\lstinline{buf[nread]} 
（对 \lstinline{nwrite} 也同样处理）。

假设对
\lstinline{piperead}
和
\lstinline{pipewrite}
的调用在两个不同的 CPU 上同时发生。
\lstinline{pipewrite}
\lineref{kernel/pipe.c:/^pipewrite/}
首先获取管道的锁，该锁保护计数、数据及其相关不变量。
\lstinline{piperead}
\lineref{kernel/pipe.c:/^piperead/}
随后也尝试获取锁，但无法获取。
它会在
\lstinline{acquire}
\lineref{kernel/spinlock.c:/^acquire/}
中自旋，等待锁。
当
\lstinline{piperead}
等待时，
\lstinline{pipewrite}
循环处理正在写入的字节
（\lstinline{addr[0..n-1]}），
依次将每个字节添加到管道中
\lineref{kernel/pipe.c:/nwrite\+\+/}。
在这个循环过程中，可能会发生缓冲区填满的情况
\lineref{kernel/pipe.c:/DOC: pipewrite-full/}。
在这种情况下，
\lstinline{pipewrite}
调用
\lstinline{wakeup}
来提醒任何正在睡眠的读者缓冲区中有数据等待，
然后在
\lstinline{&pi->nwrite}
上睡眠，等待读者从缓冲区中取出一些字节。
\lstinline{sleep}
在将
\lstinline{pipewrite}
的进程置于睡眠状态时，会释放管道锁。

\lstinline{piperead}
现在获取了管道锁并进入其临界区：
它发现
\lstinline{pi->nread}
\lstinline{!=}
\lstinline{pi->nwrite}
\lineref{kernel/pipe.c:/DOC: pipe-empty/}
（\lstinline{pipewrite}
之所以进入睡眠，是因为
\lstinline{pi->nwrite}
\lstinline{==}
\lstinline{pi->nread}
\lstinline{+}
\lstinline{PIPESIZE}
\lineref{kernel/pipe.c:/pipewrite-full/}），
因此它进入
\lstinline{for}
循环，从管道中复制数据
\lineref{kernel/pipe.c:/DOC: piperead-copy/}，
并将
\lstinline{nread}
增加复制的字节数。
现在缓冲区中有了那么多可写入的空间，因此
\lstinline{piperead}
在返回之前调用
\lstinline{wakeup}
\lineref{kernel/pipe.c:/DOC: piperead-wakeup/}
来唤醒任何正在睡眠的写入者。
\lstinline{wakeup}
找到一个在
\lstinline{&pi->nwrite}
上睡眠的进程，即正在运行
\lstinline{pipewrite}
但因缓冲区满而停止的进程。
它将该进程标记为
\indexcode{RUNNABLE}。

管道代码为读者和写者使用了不同的等待通道
（\lstinline{pi->nread}
和
\lstinline{pi->nwrite}）；
在极少数情况下，如果有大量读者和写者等待同一个管道，这可能会提高系统效率。
管道代码在循环内睡眠，该循环会检查睡眠条件；
如果有多个读者或写者，除了第一个被唤醒的进程外，其他所有进程都会发现条件不成立，并再次进入睡眠。

%% 
\section{代码：Wait、Exit 和 Kill}

请阅读 {\tt kernel/proc.c} 中函数 {\tt kwait()}、{\tt kexit()} 和 {\tt kkill()} 的代码；
它们是相应系统调用的内部实现。

%% 
\lstinline{sleep}
和
\lstinline{wakeup}
可用于多种等待场景。
一个有趣的例子，在第~\ref{CH:UNIX}~章中介绍过，是子进程的 \indexcode{exit}
与其父进程的 \indexcode{wait} 之间的交互。
在子进程死亡时，父进程可能已经在 {\tt wait} 中睡眠，或者正在做其他事情；
在后一种情况下，后续对 {\tt wait} 的调用必须观察到子进程的死亡，
这可能发生在子进程调用 {\tt exit} 很久之后。
xv6 记录子进程的死亡直到 {\tt wait} 观察到它的方式是：
{\tt exit} 将调用者置于 \indexcode{ZOMBIE} 状态，
并保持该状态直到父进程的 {\tt wait} 注意到它，
然后将子进程的状态更改为 {\tt UNUSED}，复制子进程的退出状态，
并将子进程的进程 ID 返回给父进程。
如果父进程在子进程之前退出，父进程会将子进程交给
\lstinline{init}
进程，该进程会一直调用 {\tt wait}；
因此每个子进程都有一个父进程来为其进行清理。
一个挑战是避免同时发生的父进程
\lstinline{wait}
和子进程
\lstinline{exit}，
以及同时发生的
\lstinline{exit} 和 \lstinline{exit} 之间的竞争和死锁。

\lstinline{kwait}，即 {\tt wait} 的内核实现，
首先获取
\lstinline{wait_lock}
\lineref{kernel/proc.c:/^kwait/}，
它充当条件锁，有助于确保 \lstinline{kwait} 不会错过来自退出子进程的 \lstinline{wakeup}。
然后 \lstinline{kwait} 扫描进程表。
如果它找到一个处于 \texttt{ZOMBIE} 状态的子进程，
它会释放该子进程的资源及其
\lstinline{proc} 结构，
将子进程的退出状态复制到提供给 \lstinline{wait} 的地址（如果非 0），
并返回子进程的进程 ID。
如果
\lstinline{kwait}
找到了子进程，但都没有退出，
它会调用
\lstinline{sleep}
来等待其中任何一个退出
\lineref{kernel/proc.c:/DOC: wait-sleep/}，
然后再次扫描。
\lstinline{kwait} 通常持有两个锁，
\lstinline{wait_lock} 和某个进程的 \lstinline{pp->lock}；
避免死锁的顺序是先 \lstinline{wait_lock}，然后是 \lstinline{pp->lock}。

\lstinline{kexit} \lineref{kernel/proc.c:/^kexit/} 记录退出
状态，释放一些资源，调用 \lstinline{reparent} 将所有子进程交给 \lstinline{init} 进程，
唤醒可能正在 \lstinline{wait} 中的父进程，将调用者标记为僵尸进程，
并永久让出 CPU。
\lstinline{kexit} 在此序列期间同时持有
\lstinline{wait_lock} 和 \lstinline{p->lock}。
它持有 \lstinline{wait_lock} 是因为
它是 \lstinline{wakeup(p->parent)} 的条件锁，
防止父进程在 \lstinline{wait} 中丢失唤醒。
\lstinline{kexit} 在此序列中也必须持有
\lstinline{p->lock}，
以防止在 \lstinline{wait} 中的父进程在子进程最终调用 \lstinline{swtch} 之前就看到子进程处于
\lstinline{ZOMBIE} 状态。
\lstinline{kexit} 以与 \lstinline{kwait} 相同的顺序获取这些锁以避免死锁。

\lstinline{kexit} 在将其状态设置为 \lstinline{ZOMBIE} 之前唤醒父进程，这看起来可能不正确，
但这是安全的：
尽管
\lstinline{wakeup}
可能导致父进程运行，
但父进程的
\lstinline{kwait}
循环中的检查在子进程的 \lstinline{p->lock}
被 {\tt scheduler} 释放之前无法检查子进程的状态，
因此
\lstinline{kwait}
在
\lstinline{kexit}
将其状态设置为
\lstinline{ZOMBIE}
\lineref{kernel/proc.c:/state.=.ZOMBIE/}
之前无法查看退出的进程。

虽然
\lstinline{exit}
允许进程自行终止，
但 \lstinline{kill} 系统调用
\lineref{kernel/proc.c:/^kkill/}
允许一个进程请求终止另一个进程。
让
\lstinline{kill}
直接销毁目标进程会过于复杂，因为目标进程可能正在另一个 CPU 上执行，
也许正处于对内核数据结构进行敏感更新的中间状态。
因此
\lstinline{kkill}
做得很少：它只是设置目标进程的
\indexcode{p->killed}
标志，并且如果它在睡眠中，则将其唤醒。
最终，目标进程将进入或离开内核，
此时 \lstinline{usertrap} 中的代码会检查
\lstinline{p->killed}
是否被设置（通过调用
\lstinline{killed}
\lineref{kernel/proc.c:/^killed/} 来检查），
如果被设置，则调用 \lstinline{kexit}。
如果目标进程在用户空间运行，它将在下次通过系统调用或定时器（或其他设备）中断进入内核时发现自己已被杀死。

如果目标进程正在
\lstinline{sleep}
中，
\lstinline{kkill}
对
\lstinline{wakeup}
的调用将导致目标进程从
\lstinline{sleep}
返回。
这可能很危险，因为正在等待的条件可能不成立。
然而，xv6 中对
\lstinline{sleep}
的调用总是被包装在一个
\lstinline{while}
循环中，该循环在
\lstinline{sleep}
返回后重新测试条件。
有些对
\lstinline{sleep}
的调用还会在循环中测试
\lstinline{p->killed}，
如果被设置，则放弃当前活动。
这仅在放弃是正确的做法时才会执行。
例如，管道读写代码
\lineref{kernel/pipe.c:/killed.pr/}
在杀死标志被设置时返回；最终代码将返回到 trap，trap 会再次检查 \lstinline{p->killed} 并退出。

一些 xv6 的
\lstinline{sleep}
循环不会检查
\lstinline{p->killed}，
因为代码处于多步系统调用的中间，该调用应该是原子的
（即，如果在中途放弃会出错）。
virtio 驱动程序
\lineref{kernel/virtio\_disk.c:/sleep.b/}
就是一个例子：它不检查
\lstinline{p->killed}，
因为磁盘操作可能是一组写入操作中的一个，所有这些操作都需要完成，文件系统才能保持正确状态。
一个在等待磁盘 I/O 时被杀死的进程，直到完成当前系统调用并且
\lstinline{usertrap} 看到杀死标志时才会退出。

\section{进程锁}

与每个进程关联的锁（\lstinline{p->lock}）是 xv6 中最复杂的锁。
将 \lstinline{p->lock} 视为在读取或写入以下任何 \lstinline{struct proc} 字段时都必须持有的锁，是一种简单的思考方式：
\lstinline{p->state}、
\lstinline{p->chan}、
\lstinline{p->killed}、
\lstinline{p->xstate} 和
\lstinline{p->pid}。
这些字段可能被其他进程或其他 CPU 上的调度器线程使用，因此它们自然必须受到锁的保护。

然而，\lstinline{p->lock} 的大多数使用都是为了保护 xv6 进程数据结构和算法的更高级别不变量。以下是 \lstinline{p->lock} 的完整功能列表：

% 很难知道如何以统一的方式表述这些内容。
% 讨论应该围绕不变量？保护数据？避免竞争？
% 确保原子性？避免特定风险？

% “原子”的表述通常并不能真正说明风险所在。

\begin{itemize}

\item 与 \lstinline{p->state} 一起，防止在分配新进程的 \lstinline{proc[]} 槽位时发生竞争。
% 使检查 UNUSED 与分配原子化（或类似）。

\item 在进程被创建或销毁期间将其隐藏起来不被看到。
% 使分配和销毁的所有步骤原子化。

\item 防止父进程的 \lstinline{wait} 回收一个已将其状态设置为 \lstinline{ZOMBIE} 但尚未让出 CPU 的进程。
% 将状态设置为 ZOMBIE 和让出 CPU 原子化。

\item 防止另一个 CPU 的调度器在某个进程将状态设置为 \lstinline{RUNNABLE} 之后、完成 \lstinline{swtch} 之前，决定运行这个即将让出的进程。
% 将设置 p->state 和 swtch() 原子化。

\item 确保只有一个 CPU 的调度器决定运行一个 \lstinline{RUNNABLE} 的进程。
% 将调度器检查 RUNNABLE 与实际运行进程原子化。或者将状态设置为 RUNNING。

\item 防止定时器中断在进程处于 \lstinline{swtch} 时导致其让出 CPU。
% 使 swtch（实际上是整个序列）相对于定时器中断是原子的。

\item 与条件锁一起，防止 \lstinline{wakeup} 忽略一个正在调用 \lstinline{sleep} 但尚未完成让出 CPU 的进程。
% 避免 conditionCheck+sleep 与 wakeup 之间的竞争？
% 使条件检查与让出原子化？

\item 防止 \lstinline{kill} 的目标进程在 \lstinline{kkill} 检查 \lstinline{p->pid} 和设置 \lstinline{p->killed} 之间退出并可能被重新分配。
% 使 pid 检查与设置 killed 原子化。

\item 使 \lstinline{kkill} 对 \lstinline{p->state} 的检查和写入原子化。

\end{itemize}

\lstinline{p->parent} 字段受全局锁
\lstinline{wait_lock} 的保护，而不是 \lstinline{p->lock}。
只有进程的父进程会修改 \lstinline{p->parent}，
尽管该字段会被进程自身以及搜索其子进程的其他进程读取。
\lstinline{wait_lock} 的目的是当 \lstinline{wait} 睡眠等待任何子进程退出时充当条件锁。
退出的子进程在将其状态设置为 \lstinline{ZOMBIE}、唤醒其父进程并让出 CPU 之前，
会一直持有 \lstinline{wait_lock} 或 \lstinline{p->lock}。
\lstinline{wait_lock} 还序列化父进程和子进程并发的 \lstinline{exit}，
从而保证 \lstinline{init} 进程（它继承子进程）从其 \lstinline{wait} 中被唤醒。
\lstinline{wait_lock} 是一个全局锁，而不是每个父进程的每进程锁，
因为一个进程在获取它之前，无法知道其父进程是谁。

%% 
\section{现实世界}
%% 

\lstinline{sleep}
和
\lstinline{wakeup}
是一种简单而有效的同步方法，
但还有许多其他的方法；
信号量~\cite{dijkstra65} 就是一个例子。
所有方法中的第一个挑战是避免
我们在本章开头看到的“丢失唤醒”问题。
最初的 Unix 内核的
\lstinline{sleep}
只是简单地禁用中断，
这在 Unix 运行在单 CPU 系统上时是足够的。
由于 xv6 运行在多处理器上，
它为
\lstinline{sleep}
添加了一个显式的锁。
FreeBSD 的
\lstinline{msleep}
采用了相同的方法。
Plan 9 的
\lstinline{sleep}
使用了一个回调函数，该函数在即将进入睡眠之前持有调度锁的情况下运行；
该函数作为对睡眠条件的最后一刻检查，以避免丢失唤醒。
Linux 内核的
\lstinline{sleep}
使用一个显式的进程队列，称为等待队列，而不是等待通道；该队列有自己的内部锁。

在 \lstinline{wakeup} 中扫描整个进程表是低效的。
一个更好的解决方案是用一个数据结构替换
\lstinline{sleep}
和
\lstinline{wakeup}
中的
\lstinline{chan}，
该数据结构保存一个在该结构上睡眠的进程列表，
例如 Linux 的等待队列。
Plan 9 的
\lstinline{sleep}
和
\lstinline{wakeup}
称这种结构为会合点。
许多线程库将相同的结构称为条件变量；
在这种情况下，操作
\lstinline{sleep}
和
\lstinline{wakeup}
被称为
\lstinline{wait}
和
\lstinline{signal}。
所有这些机制都具有相同的风格：睡眠条件受某种锁的保护，该锁在睡眠期间原子地释放。

xv6 的
\lstinline{wakeup}
会唤醒在特定等待通道上等待的所有进程。
如果其中有多个进程，它们都会尝试获取条件锁并重新检查条件；
在许多情况下，只有一个进程能够执行有用的操作（例如，读取管道中等待的所有数据）。
其余的会发现条件不再成立，并重新进入睡眠状态；
唤醒它们浪费了 CPU 时间。
因此，大多数条件变量设计提供了两个原语：
\lstinline{signal}，
它唤醒等待条件变量的进程之一，以及
\lstinline{broadcast}，
它会唤醒所有等待的进程。

强制杀死进程会带来一些问题。
例如，一个被杀死的进程可能深度处于内核睡眠中，解开其堆栈需要小心，
因为调用堆栈上的每个函数可能都需要进行一些清理。
某些语言通过提供异常机制来提供帮助，但 C 语言不行。
此外，还有其他事件可能导致正在睡眠的进程被唤醒，即使它正在等待的事件尚未发生。
例如，当一个 Unix 进程处于睡眠状态时，另一个进程可能会向其发送一个
\indexcode{signal}。
在这种情况下，进程将从被中断的系统调用返回，返回值为 -1，错误码设置为 EINTR。
应用程序可以检查这些值并决定如何处理。
Xv6 不支持信号，因此不会出现这种复杂性。

Xv6 对 \lstinline{kill} 的支持并不完全令人满意：
有些睡眠循环可能应该检查 \lstinline{p->killed}。
一个相关的问题是，即使对于检查 \lstinline{p->killed} 的睡眠循环，
在 \lstinline{sleep} 和 \lstinline{kill} 之间也存在竞争；
后者可能刚好在受害者的循环检查完 \lstinline{p->killed} 之后、但在其调用 \lstinline{sleep} 之前设置 \lstinline{p->killed} 并尝试唤醒受害者。
如果发生此问题，受害者只有在它等待的条件发生时才会注意到 \lstinline{p->killed}。
这可能会延迟很久甚至永远不会发生（例如，如果受害者在等待控制台输入，而用户没有输入任何字符）。

%% 
\section{练习}
%% 

\begin{enumerate}

\item 在 xv6 中实现计数信号量。
选择 xv6 中几个使用 sleep 和 wakeup 的地方，用信号量替换它们。
评判结果。

\item 你能实现一个 {\tt sleep()} 的变体，它只接受一个参数（通道），而不需要锁参数吗？

\item 修复上述 \lstinline{kill} 和 \lstinline{sleep} 之间的竞争条件，
使得在受害者的睡眠循环检查完 \lstinline{p->killed} 之后、但在其调用 \lstinline{sleep} 之前发生的 \lstinline{kill} 能够导致受害者放弃当前系统调用。
% 答案：一种解决方案是在 sleep 中，在将进程状态设置为 sleep 之前检查 p->killed 是否被设置。

\item 设计一个方案，使每个睡眠循环都检查 \lstinline{p->killed}，
这样，例如，virtio 驱动程序中的进程如果被另一个进程杀死，就可以快速从 while 循环中返回。
% 答案：这很困难。现代 Unix 系统使用 setjmp 和 longjmp 以及非常仔细的编程来清理被中断的系统调用可能已构建的任何部分状态。

\end{enumerate}

  \caption{重叠锁以避免丢失唤醒}
    \label{fig:overlap}
\end{figure}

为什么 \lstinline{sleep} 和 \lstinline{wakeup} 的锁规则能确保即将进入睡眠的进程不会错过并发的唤醒？
即将进入睡眠的进程从{\it 在}检查条件之前，到{\it 在}将其自身标记为 \texttt{SLEEPING} 之后，
始终持有条件锁或其自身的 \lstinline{p->lock} 或两者都持有；
见图~\ref{fig:overlap}。
调用 \texttt{wakeup} 的进程需要获取{\it 两个}锁。
唤醒者可能先获取锁，这意味着它会在消费者线程检查条件之前使条件成立，
这样消费者线程就不需要调用 {\tt sleep()}；
或者唤醒者的 {\tt acquire()} 可能需要等待，直到消费者线程完全完成进入睡眠并释放锁，
此时唤醒者将看到消费者线程已被标记为 {\tt SLEEPING} 并将其唤醒。

有时多个进程在同一个通道上睡眠；例如，多个进程从一个管道读取数据。
一次对 \lstinline{wakeup} 的调用将唤醒它们全部。
其中一个将首先运行，并获取 \lstinline{sleep} 调用时所持有的锁，
并且（在管道的例子中）读取所有等待的数据。
其他进程会发现，尽管被唤醒了，但没有数据可读。
从它们的角度来看，这次唤醒是“虚假的”，它们必须再次进入睡眠。
因此，\lstinline{sleep} 总是在一个重新检查条件的循环内被调用，如上面的 \lstinline{P} 所示。

如果两个使用 sleep/wakeup 的地方意外选择了相同的通道，也不会造成伤害：
它们会遇到虚假唤醒，但如上所述的循环可以容忍这个问题。
sleep/wakeup 的魅力很大程度上在于它既轻量级（无需创建特殊的数据结构作为等待通道），
又提供了一层间接性（调用者无需知道他们正在与哪个特定进程交互）。

%% 
\section{代码：管道}
%% 
xv6 的管道是使用 \lstinline{sleep} 和 \lstinline{wakeup} 来同步生产者和消费者的代码示例。
我们在第~\ref{CH:UNIX}~章中了解了管道的接口：
写入管道一端的数据被复制到内核缓冲区，然后可以从另一端读取。
后续章节将研究围绕管道的文件描述符支持，现在让我们看看
\indexcode{pipewrite}
和
\indexcode{piperead}
的实现。

每个管道由一个
\indexcode{struct pipe}
表示，其中包含一个
\lstinline{lock}
和一个
\lstinline{data}
缓冲区。
字段
\lstinline{nread}
和
\lstinline{nwrite}
统计从缓冲区读取和写入缓冲区的总字节数。
缓冲区是环绕的：
在
\lstinline{buf[PIPESIZE-1]}
之后写入的下一个字节是
\lstinline{buf[0]}。
计数不会环绕。
这个约定使得实现能够区分缓冲区满
（\lstinline{nwrite}
\lstinline{==}
\lstinline{nread+PIPESIZE}）
和缓冲区空
（\lstinline{nwrite}
\lstinline{==}
\lstinline{nread}），
但这意味着对缓冲区的索引必须使用
\lstinline{buf[nread}
\lstinline{%}
\lstinline{PIPESIZE]}
而不是简单的
\lstinline{buf[nread]} 
（对 \lstinline{nwrite} 也同样处理）。

假设对
\lstinline{piperead}
和
\lstinline{pipewrite}
的调用在两个不同的 CPU 上同时发生。
\lstinline{pipewrite}
\lineref{kernel/pipe.c:/^pipewrite/}
首先获取管道的锁，该锁保护计数、数据及其相关不变量。
\lstinline{piperead}
\lineref{kernel/pipe.c:/^piperead/}
随后也尝试获取锁，但无法获取。
它会在
\lstinline{acquire}
\lineref{kernel/spinlock.c:/^acquire/}
中自旋，等待锁。
当
\lstinline{piperead}
等待时，
\lstinline{pipewrite}
循环处理正在写入的字节
（\lstinline{addr[0..n-1]}），
依次将每个字节添加到管道中
\lineref{kernel/pipe.c:/nwrite\+\+/}。
在这个循环过程中，可能会发生缓冲区填满的情况
\lineref{kernel/pipe.c:/DOC: pipewrite-full/}。
在这种情况下，
\lstinline{pipewrite}
调用
\lstinline{wakeup}
来提醒任何正在睡眠的读者缓冲区中有数据等待，
然后在
\lstinline{&pi->nwrite}
上睡眠，等待读者从缓冲区中取出一些字节。
\lstinline{sleep}
在将
\lstinline{pipewrite}
的进程置于睡眠状态时，会释放管道锁。

\lstinline{piperead}
现在获取了管道锁并进入其临界区：
它发现
\lstinline{pi->nread}
\lstinline{!=}
\lstinline{pi->nwrite}
\lineref{kernel/pipe.c:/DOC: pipe-empty/}
（\lstinline{pipewrite}
之所以进入睡眠，是因为
\lstinline{pi->nwrite}
\lstinline{==}
\lstinline{pi->nread}
\lstinline{+}
\lstinline{PIPESIZE}
\lineref{kernel/pipe.c:/pipewrite-full/}），
因此它进入
\lstinline{for}
循环，从管道中复制数据
\lineref{kernel/pipe.c:/DOC: piperead-copy/}，
并将
\lstinline{nread}
增加复制的字节数。
现在缓冲区中有了那么多可写入的空间，因此
\lstinline{piperead}
在返回之前调用
\lstinline{wakeup}
\lineref{kernel/pipe.c:/DOC: piperead-wakeup/}
来唤醒任何正在睡眠的写入者。
\lstinline{wakeup}
找到一个在
\lstinline{&pi->nwrite}
上睡眠的进程，即正在运行
\lstinline{pipewrite}
但因缓冲区满而停止的进程。
它将该进程标记为
\indexcode{RUNNABLE}。

管道代码为读者和写者使用了不同的等待通道
（\lstinline{pi->nread}
和
\lstinline{pi->nwrite}）；
在极少数情况下，如果有大量读者和写者等待同一个管道，这可能会提高系统效率。
管道代码在循环内睡眠，该循环会检查睡眠条件；
如果有多个读者或写者，除了第一个被唤醒的进程外，其他所有进程都会发现条件不成立，并再次进入睡眠。

%% 
\section{代码：Wait、Exit 和 Kill}

请阅读 {\tt kernel/proc.c} 中函数 {\tt kwait()}、{\tt kexit()} 和 {\tt kkill()} 的代码；
它们是相应系统调用的内部实现。

%% 
\lstinline{sleep}
和
\lstinline{wakeup}
可用于多种等待场景。
一个有趣的例子，在第~\ref{CH:UNIX}~章中介绍过，是子进程的 \indexcode{exit}
与其父进程的 \indexcode{wait} 之间的交互。
在子进程死亡时，父进程可能已经在 {\tt wait} 中睡眠，或者正在做其他事情；
在后一种情况下，后续对 {\tt wait} 的调用必须观察到子进程的死亡，
这可能发生在子进程调用 {\tt exit} 很久之后。
xv6 记录子进程的死亡直到 {\tt wait} 观察到它的方式是：
{\tt exit} 将调用者置于 \indexcode{ZOMBIE} 状态，
并保持该状态直到父进程的 {\tt wait} 注意到它，
然后将子进程的状态更改为 {\tt UNUSED}，复制子进程的退出状态，
并将子进程的进程 ID 返回给父进程。
如果父进程在子进程之前退出，父进程会将子进程交给
\lstinline{init}
进程，该进程会一直调用 {\tt wait}；
因此每个子进程都有一个父进程来为其进行清理。
一个挑战是避免同时发生的父进程
\lstinline{wait}
和子进程
\lstinline{exit}，
以及同时发生的
\lstinline{exit} 和 \lstinline{exit} 之间的竞争和死锁。

\lstinline{kwait}，即 {\tt wait} 的内核实现，
首先获取
\lstinline{wait_lock}
\lineref{kernel/proc.c:/^kwait/}，
它充当条件锁，有助于确保 \lstinline{kwait} 不会错过来自退出子进程的 \lstinline{wakeup}。
然后 \lstinline{kwait} 扫描进程表。
如果它找到一个处于 \texttt{ZOMBIE} 状态的子进程，
它会释放该子进程的资源及其
\lstinline{proc} 结构，
将子进程的退出状态复制到提供给 \lstinline{wait} 的地址（如果非 0），
并返回子进程的进程 ID。
如果
\lstinline{kwait}
找到了子进程，但都没有退出，
它会调用
\lstinline{sleep}
来等待其中任何一个退出
\lineref{kernel/proc.c:/DOC: wait-sleep/}，
然后再次扫描。
\lstinline{kwait} 通常持有两个锁，
\lstinline{wait_lock} 和某个进程的 \lstinline{pp->lock}；
避免死锁的顺序是先 \lstinline{wait_lock}，然后是 \lstinline{pp->lock}。

\lstinline{kexit} \lineref{kernel/proc.c:/^kexit/} 记录退出
状态，释放一些资源，调用 \lstinline{reparent} 将所有子进程交给 \lstinline{init} 进程，
唤醒可能正在 \lstinline{wait} 中的父进程，将调用者标记为僵尸进程，
并永久让出 CPU。
\lstinline{kexit} 在此序列期间同时持有
\lstinline{wait_lock} 和 \lstinline{p->lock}。
它持有 \lstinline{wait_lock} 是因为
它是 \lstinline{wakeup(p->parent)} 的条件锁，
防止父进程在 \lstinline{wait} 中丢失唤醒。
\lstinline{kexit} 在此序列中也必须持有
\lstinline{p->lock}，
以防止在 \lstinline{wait} 中的父进程在子进程最终调用 \lstinline{swtch} 之前就看到子进程处于
\lstinline{ZOMBIE} 状态。
\lstinline{kexit} 以与 \lstinline{kwait} 相同的顺序获取这些锁以避免死锁。

\lstinline{kexit} 在将其状态设置为 \lstinline{ZOMBIE} 之前唤醒父进程，这看起来可能不正确，
但这是安全的：
尽管
\lstinline{wakeup}
可能导致父进程运行，
但父进程的
\lstinline{kwait}
循环中的检查在子进程的 \lstinline{p->lock}
被 {\tt scheduler} 释放之前无法检查子进程的状态，
因此
\lstinline{kwait}
在
\lstinline{kexit}
将其状态设置为
\lstinline{ZOMBIE}
\lineref{kernel/proc.c:/state.=.ZOMBIE/}
之前无法查看退出的进程。

虽然
\lstinline{exit}
允许进程自行终止，
但 \lstinline{kill} 系统调用
\lineref{kernel/proc.c:/^kkill/}
允许一个进程请求终止另一个进程。
让
\lstinline{kill}
直接销毁目标进程会过于复杂，因为目标进程可能正在另一个 CPU 上执行，
也许正处于对内核数据结构进行敏感更新的中间状态。
因此
\lstinline{kkill}
做得很少：它只是设置目标进程的
\indexcode{p->killed}
标志，并且如果它在睡眠中，则将其唤醒。
最终，目标进程将进入或离开内核，
此时 \lstinline{usertrap} 中的代码会检查
\lstinline{p->killed}
是否被设置（通过调用
\lstinline{killed}
\lineref{kernel/proc.c:/^killed/} 来检查），
如果被设置，则调用 \lstinline{kexit}。
如果目标进程在用户空间运行，它将在下次通过系统调用或定时器（或其他设备）中断进入内核时发现自己已被杀死。

如果目标进程正在
\lstinline{sleep}
中，
\lstinline{kkill}
对
\lstinline{wakeup}
的调用将导致目标进程从
\lstinline{sleep}
返回。
这可能很危险，因为正在等待的条件可能不成立。
然而，xv6 中对
\lstinline{sleep}
的调用总是被包装在一个
\lstinline{while}
循环中，该循环在
\lstinline{sleep}
返回后重新测试条件。
有些对
\lstinline{sleep}
的调用还会在循环中测试
\lstinline{p->killed}，
如果被设置，则放弃当前活动。
这仅在放弃是正确的做法时才会执行。
例如，管道读写代码
\lineref{kernel/pipe.c:/killed.pr/}
在杀死标志被设置时返回；最终代码将返回到 trap，trap 会再次检查 \lstinline{p->killed} 并退出。

一些 xv6 的
\lstinline{sleep}
循环不会检查
\lstinline{p->killed}，
因为代码处于多步系统调用的中间，该调用应该是原子的
（即，如果在中途放弃会出错）。
virtio 驱动程序
\lineref{kernel/virtio\_disk.c:/sleep.b/}
就是一个例子：它不检查
\lstinline{p->killed}，
因为磁盘操作可能是一组写入操作中的一个，所有这些操作都需要完成，文件系统才能保持正确状态。
一个在等待磁盘 I/O 时被杀死的进程，直到完成当前系统调用并且
\lstinline{usertrap} 看到杀死标志时才会退出。

\section{进程锁}

与每个进程关联的锁（\lstinline{p->lock}）是 xv6 中最复杂的锁。
将 \lstinline{p->lock} 视为在读取或写入以下任何 \lstinline{struct proc} 字段时都必须持有的锁，是一种简单的思考方式：
\lstinline{p->state}、
\lstinline{p->chan}、
\lstinline{p->killed}、
\lstinline{p->xstate} 和
\lstinline{p->pid}。
这些字段可能被其他进程或其他 CPU 上的调度器线程使用，因此它们自然必须受到锁的保护。

然而，\lstinline{p->lock} 的大多数使用都是为了保护 xv6 进程数据结构和算法的更高级别不变量。以下是 \lstinline{p->lock} 的完整功能列表：

% 很难知道如何以统一的方式表述这些内容。
% 讨论应该围绕不变量？保护数据？避免竞争？
% 确保原子性？避免特定风险？

% “原子”的表述通常并不能真正说明风险所在。

\begin{itemize}

\item 与 \lstinline{p->state} 一起，防止在分配新进程的 \lstinline{proc[]} 槽位时发生竞争。
% 使检查 UNUSED 与分配原子化（或类似）。

\item 在进程被创建或销毁期间将其隐藏起来不被看到。
% 使分配和销毁的所有步骤原子化。

\item 防止父进程的 \lstinline{wait} 回收一个已将其状态设置为 \lstinline{ZOMBIE} 但尚未让出 CPU 的进程。
% 将状态设置为 ZOMBIE 和让出 CPU 原子化。

\item 防止另一个 CPU 的调度器在某个进程将状态设置为 \lstinline{RUNNABLE} 之后、完成 \lstinline{swtch} 之前，决定运行这个即将让出的进程。
% 将设置 p->state 和 swtch() 原子化。

\item 确保只有一个 CPU 的调度器决定运行一个 \lstinline{RUNNABLE} 的进程。
% 将调度器检查 RUNNABLE 与实际运行进程原子化。或者将状态设置为 RUNNING。

\item 防止定时器中断在进程处于 \lstinline{swtch} 时导致其让出 CPU。
% 使 swtch（实际上是整个序列）相对于定时器中断是原子的。

\item 与条件锁一起，防止 \lstinline{wakeup} 忽略一个正在调用 \lstinline{sleep} 但尚未完成让出 CPU 的进程。
% 避免 conditionCheck+sleep 与 wakeup 之间的竞争？
% 使条件检查与让出原子化？

\item 防止 \lstinline{kill} 的目标进程在 \lstinline{kkill} 检查 \lstinline{p->pid} 和设置 \lstinline{p->killed} 之间退出并可能被重新分配。
% 使 pid 检查与设置 killed 原子化。

\item 使 \lstinline{kkill} 对 \lstinline{p->state} 的检查和写入原子化。

\end{itemize}

\lstinline{p->parent} 字段受全局锁
\lstinline{wait_lock} 的保护，而不是 \lstinline{p->lock}。
只有进程的父进程会修改 \lstinline{p->parent}，
尽管该字段会被进程自身以及搜索其子进程的其他进程读取。
\lstinline{wait_lock} 的目的是当 \lstinline{wait} 睡眠等待任何子进程退出时充当条件锁。
退出的子进程在将其状态设置为 \lstinline{ZOMBIE}、唤醒其父进程并让出 CPU 之前，
会一直持有 \lstinline{wait_lock} 或 \lstinline{p->lock}。
\lstinline{wait_lock} 还序列化父进程和子进程并发的 \lstinline{exit}，
从而保证 \lstinline{init} 进程（它继承子进程）从其 \lstinline{wait} 中被唤醒。
\lstinline{wait_lock} 是一个全局锁，而不是每个父进程的每进程锁，
因为一个进程在获取它之前，无法知道其父进程是谁。

%% 
\section{现实世界}
%% 

\lstinline{sleep}
和
\lstinline{wakeup}
是一种简单而有效的同步方法，
但还有许多其他的方法；
信号量~\cite{dijkstra65} 就是一个例子。
所有方法中的第一个挑战是避免
我们在本章开头看到的“丢失唤醒”问题。
最初的 Unix 内核的
\lstinline{sleep}
只是简单地禁用中断，
这在 Unix 运行在单 CPU 系统上时是足够的。
由于 xv6 运行在多处理器上，
它为
\lstinline{sleep}
添加了一个显式的锁。
FreeBSD 的
\lstinline{msleep}
采用了相同的方法。
Plan 9 的
\lstinline{sleep}
使用了一个回调函数，该函数在即将进入睡眠之前持有调度锁的情况下运行；
该函数作为对睡眠条件的最后一刻检查，以避免丢失唤醒。
Linux 内核的
\lstinline{sleep}
使用一个显式的进程队列，称为等待队列，而不是等待通道；该队列有自己的内部锁。

在 \lstinline{wakeup} 中扫描整个进程表是低效的。
一个更好的解决方案是用一个数据结构替换
\lstinline{sleep}
和
\lstinline{wakeup}
中的
\lstinline{chan}，
该数据结构保存一个在该结构上睡眠的进程列表，
例如 Linux 的等待队列。
Plan 9 的
\lstinline{sleep}
和
\lstinline{wakeup}
称这种结构为会合点。
许多线程库将相同的结构称为条件变量；
在这种情况下，操作
\lstinline{sleep}
和
\lstinline{wakeup}
被称为
\lstinline{wait}
和
\lstinline{signal}。
所有这些机制都具有相同的风格：睡眠条件受某种锁的保护，该锁在睡眠期间原子地释放。

xv6 的
\lstinline{wakeup}
会唤醒在特定等待通道上等待的所有进程。
如果其中有多个进程，它们都会尝试获取条件锁并重新检查条件；
在许多情况下，只有一个进程能够执行有用的操作（例如，读取管道中等待的所有数据）。
其余的会发现条件不再成立，并重新进入睡眠状态；
唤醒它们浪费了 CPU 时间。
因此，大多数条件变量设计提供了两个原语：
\lstinline{signal}，
它唤醒等待条件变量的进程之一，以及
\lstinline{broadcast}，
它会唤醒所有等待的进程。

强制杀死进程会带来一些问题。
例如，一个被杀死的进程可能深度处于内核睡眠中，解开其堆栈需要小心，
因为调用堆栈上的每个函数可能都需要进行一些清理。
某些语言通过提供异常机制来提供帮助，但 C 语言不行。
此外，还有其他事件可能导致正在睡眠的进程被唤醒，即使它正在等待的事件尚未发生。
例如，当一个 Unix 进程处于睡眠状态时，另一个进程可能会向其发送一个
\indexcode{signal}。
在这种情况下，进程将从被中断的系统调用返回，返回值为 -1，错误码设置为 EINTR。
应用程序可以检查这些值并决定如何处理。
Xv6 不支持信号，因此不会出现这种复杂性。

Xv6 对 \lstinline{kill} 的支持并不完全令人满意：
有些睡眠循环可能应该检查 \lstinline{p->killed}。
一个相关的问题是，即使对于检查 \lstinline{p->killed} 的睡眠循环，
在 \lstinline{sleep} 和 \lstinline{kill} 之间也存在竞争；
后者可能刚好在受害者的循环检查完 \lstinline{p->killed} 之后、但在其调用 \lstinline{sleep} 之前设置 \lstinline{p->killed} 并尝试唤醒受害者。
如果发生此问题，受害者只有在它等待的条件发生时才会注意到 \lstinline{p->killed}。
这可能会延迟很久甚至永远不会发生（例如，如果受害者在等待控制台输入，而用户没有输入任何字符）。

%% 
\section{练习}
%% 

\begin{enumerate}

\item 在 xv6 中实现计数信号量。
选择 xv6 中几个使用 sleep 和 wakeup 的地方，用信号量替换它们。
评判结果。

\item 你能实现一个 {\tt sleep()} 的变体，它只接受一个参数（通道），而不需要锁参数吗？

\item 修复上述 \lstinline{kill} 和 \lstinline{sleep} 之间的竞争条件，
使得在受害者的睡眠循环检查完 \lstinline{p->killed} 之后、但在其调用 \lstinline{sleep} 之前发生的 \lstinline{kill} 能够导致受害者放弃当前系统调用。
% 答案：一种解决方案是在 sleep 中，在将进程状态设置为 sleep 之前检查 p->killed 是否被设置。

\item 设计一个方案，使每个睡眠循环都检查 \lstinline{p->killed}，
这样，例如，virtio 驱动程序中的进程如果被另一个进程杀死，就可以快速从 while 循环中返回。
% 答案：这很困难。现代 Unix 系统使用 setjmp 和 longjmp 以及非常仔细的编程来清理被中断的系统调用可能已构建的任何部分状态。

\end{enumerate}

  \caption{重叠锁以避免丢失唤醒}
    \label{fig:overlap}
\end{figure}

为什么 \lstinline{sleep} 和 \lstinline{wakeup} 的锁规则能确保即将进入睡眠的进程不会错过并发的唤醒？
即将进入睡眠的进程从{\it 在}检查条件之前，到{\it 在}将其自身标记为 \texttt{SLEEPING} 之后，
始终持有条件锁或其自身的 \lstinline{p->lock} 或两者都持有；
见图~\ref{fig:overlap}。
调用 \texttt{wakeup} 的进程需要获取{\it 两个}锁。
唤醒者可能先获取锁，这意味着它会在消费者线程检查条件之前使条件成立，
这样消费者线程就不需要调用 {\tt sleep()}；
或者唤醒者的 {\tt acquire()} 可能需要等待，直到消费者线程完全完成进入睡眠并释放锁，
此时唤醒者将看到消费者线程已被标记为 {\tt SLEEPING} 并将其唤醒。

有时多个进程在同一个通道上睡眠；例如，多个进程从一个管道读取数据。
一次对 \lstinline{wakeup} 的调用将唤醒它们全部。
其中一个将首先运行，并获取 \lstinline{sleep} 调用时所持有的锁，
并且（在管道的例子中）读取所有等待的数据。
其他进程会发现，尽管被唤醒了，但没有数据可读。
从它们的角度来看，这次唤醒是“虚假的”，它们必须再次进入睡眠。
因此，\lstinline{sleep} 总是在一个重新检查条件的循环内被调用，如上面的 \lstinline{P} 所示。

如果两个使用 sleep/wakeup 的地方意外选择了相同的通道，也不会造成伤害：
它们会遇到虚假唤醒，但如上所述的循环可以容忍这个问题。
sleep/wakeup 的魅力很大程度上在于它既轻量级（无需创建特殊的数据结构作为等待通道），
又提供了一层间接性（调用者无需知道他们正在与哪个特定进程交互）。

%% 
\section{代码：管道}
%% 
xv6 的管道是使用 \lstinline{sleep} 和 \lstinline{wakeup} 来同步生产者和消费者的代码示例。
我们在第~\ref{CH:UNIX}~章中了解了管道的接口：
写入管道一端的数据被复制到内核缓冲区，然后可以从另一端读取。
后续章节将研究围绕管道的文件描述符支持，现在让我们看看
\indexcode{pipewrite}
和
\indexcode{piperead}
的实现。

每个管道由一个
\indexcode{struct pipe}
表示，其中包含一个
\lstinline{lock}
和一个
\lstinline{data}
缓冲区。
字段
\lstinline{nread}
和
\lstinline{nwrite}
统计从缓冲区读取和写入缓冲区的总字节数。
缓冲区是环绕的：
在
\lstinline{buf[PIPESIZE-1]}
之后写入的下一个字节是
\lstinline{buf[0]}。
计数不会环绕。
这个约定使得实现能够区分缓冲区满
（\lstinline{nwrite}
\lstinline{==}
\lstinline{nread+PIPESIZE}）
和缓冲区空
（\lstinline{nwrite}
\lstinline{==}
\lstinline{nread}），
但这意味着对缓冲区的索引必须使用
\lstinline{buf[nread}
\lstinline{%}
\lstinline{PIPESIZE]}
而不是简单的
\lstinline{buf[nread]} 
（对 \lstinline{nwrite} 也同样处理）。

假设对
\lstinline{piperead}
和
\lstinline{pipewrite}
的调用在两个不同的 CPU 上同时发生。
\lstinline{pipewrite}
\lineref{kernel/pipe.c:/^pipewrite/}
首先获取管道的锁，该锁保护计数、数据及其相关不变量。
\lstinline{piperead}
\lineref{kernel/pipe.c:/^piperead/}
随后也尝试获取锁，但无法获取。
它会在
\lstinline{acquire}
\lineref{kernel/spinlock.c:/^acquire/}
中自旋，等待锁。
当
\lstinline{piperead}
等待时，
\lstinline{pipewrite}
循环处理正在写入的字节
（\lstinline{addr[0..n-1]}），
依次将每个字节添加到管道中
\lineref{kernel/pipe.c:/nwrite\+\+/}。
在这个循环过程中，可能会发生缓冲区填满的情况
\lineref{kernel/pipe.c:/DOC: pipewrite-full/}。
在这种情况下，
\lstinline{pipewrite}
调用
\lstinline{wakeup}
来提醒任何正在睡眠的读者缓冲区中有数据等待，
然后在
\lstinline{&pi->nwrite}
上睡眠，等待读者从缓冲区中取出一些字节。
\lstinline{sleep}
在将
\lstinline{pipewrite}
的进程置于睡眠状态时，会释放管道锁。

\lstinline{piperead}
现在获取了管道锁并进入其临界区：
它发现
\lstinline{pi->nread}
\lstinline{!=}
\lstinline{pi->nwrite}
\lineref{kernel/pipe.c:/DOC: pipe-empty/}
（\lstinline{pipewrite}
之所以进入睡眠，是因为
\lstinline{pi->nwrite}
\lstinline{==}
\lstinline{pi->nread}
\lstinline{+}
\lstinline{PIPESIZE}
\lineref{kernel/pipe.c:/pipewrite-full/}），
因此它进入
\lstinline{for}
循环，从管道中复制数据
\lineref{kernel/pipe.c:/DOC: piperead-copy/}，
并将
\lstinline{nread}
增加复制的字节数。
现在缓冲区中有了那么多可写入的空间，因此
\lstinline{piperead}
在返回之前调用
\lstinline{wakeup}
\lineref{kernel/pipe.c:/DOC: piperead-wakeup/}
来唤醒任何正在睡眠的写入者。
\lstinline{wakeup}
找到一个在
\lstinline{&pi->nwrite}
上睡眠的进程，即正在运行
\lstinline{pipewrite}
但因缓冲区满而停止的进程。
它将该进程标记为
\indexcode{RUNNABLE}。

管道代码为读者和写者使用了不同的等待通道
（\lstinline{pi->nread}
和
\lstinline{pi->nwrite}）；
在极少数情况下，如果有大量读者和写者等待同一个管道，这可能会提高系统效率。
管道代码在循环内睡眠，该循环会检查睡眠条件；
如果有多个读者或写者，除了第一个被唤醒的进程外，其他所有进程都会发现条件不成立，并再次进入睡眠。

%% 
\section{代码：Wait、Exit 和 Kill}

请阅读 {\tt kernel/proc.c} 中函数 {\tt kwait()}、{\tt kexit()} 和 {\tt kkill()} 的代码；
它们是相应系统调用的内部实现。

%% 
\lstinline{sleep}
和
\lstinline{wakeup}
可用于多种等待场景。
一个有趣的例子，在第~\ref{CH:UNIX}~章中介绍过，是子进程的 \indexcode{exit}
与其父进程的 \indexcode{wait} 之间的交互。
在子进程死亡时，父进程可能已经在 {\tt wait} 中睡眠，或者正在做其他事情；
在后一种情况下，后续对 {\tt wait} 的调用必须观察到子进程的死亡，
这可能发生在子进程调用 {\tt exit} 很久之后。
xv6 记录子进程的死亡直到 {\tt wait} 观察到它的方式是：
{\tt exit} 将调用者置于 \indexcode{ZOMBIE} 状态，
并保持该状态直到父进程的 {\tt wait} 注意到它，
然后将子进程的状态更改为 {\tt UNUSED}，复制子进程的退出状态，
并将子进程的进程 ID 返回给父进程。
如果父进程在子进程之前退出，父进程会将子进程交给
\lstinline{init}
进程，该进程会一直调用 {\tt wait}；
因此每个子进程都有一个父进程来为其进行清理。
一个挑战是避免同时发生的父进程
\lstinline{wait}
和子进程
\lstinline{exit}，
以及同时发生的
\lstinline{exit} 和 \lstinline{exit} 之间的竞争和死锁。

\lstinline{kwait}，即 {\tt wait} 的内核实现，
首先获取
\lstinline{wait_lock}
\lineref{kernel/proc.c:/^kwait/}，
它充当条件锁，有助于确保 \lstinline{kwait} 不会错过来自退出子进程的 \lstinline{wakeup}。
然后 \lstinline{kwait} 扫描进程表。
如果它找到一个处于 \texttt{ZOMBIE} 状态的子进程，
它会释放该子进程的资源及其
\lstinline{proc} 结构，
将子进程的退出状态复制到提供给 \lstinline{wait} 的地址（如果非 0），
并返回子进程的进程 ID。
如果
\lstinline{kwait}
找到了子进程，但都没有退出，
它会调用
\lstinline{sleep}
来等待其中任何一个退出
\lineref{kernel/proc.c:/DOC: wait-sleep/}，
然后再次扫描。
\lstinline{kwait} 通常持有两个锁，
\lstinline{wait_lock} 和某个进程的 \lstinline{pp->lock}；
避免死锁的顺序是先 \lstinline{wait_lock}，然后是 \lstinline{pp->lock}。

\lstinline{kexit} \lineref{kernel/proc.c:/^kexit/} 记录退出
状态，释放一些资源，调用 \lstinline{reparent} 将所有子进程交给 \lstinline{init} 进程，
唤醒可能正在 \lstinline{wait} 中的父进程，将调用者标记为僵尸进程，
并永久让出 CPU。
\lstinline{kexit} 在此序列期间同时持有
\lstinline{wait_lock} 和 \lstinline{p->lock}。
它持有 \lstinline{wait_lock} 是因为
它是 \lstinline{wakeup(p->parent)} 的条件锁，
防止父进程在 \lstinline{wait} 中丢失唤醒。
\lstinline{kexit} 在此序列中也必须持有
\lstinline{p->lock}，
以防止在 \lstinline{wait} 中的父进程在子进程最终调用 \lstinline{swtch} 之前就看到子进程处于
\lstinline{ZOMBIE} 状态。
\lstinline{kexit} 以与 \lstinline{kwait} 相同的顺序获取这些锁以避免死锁。

\lstinline{kexit} 在将其状态设置为 \lstinline{ZOMBIE} 之前唤醒父进程，这看起来可能不正确，
但这是安全的：
尽管
\lstinline{wakeup}
可能导致父进程运行，
但父进程的
\lstinline{kwait}
循环中的检查在子进程的 \lstinline{p->lock}
被 {\tt scheduler} 释放之前无法检查子进程的状态，
因此
\lstinline{kwait}
在
\lstinline{kexit}
将其状态设置为
\lstinline{ZOMBIE}
\lineref{kernel/proc.c:/state.=.ZOMBIE/}
之前无法查看退出的进程。

虽然
\lstinline{exit}
允许进程自行终止，
但 \lstinline{kill} 系统调用
\lineref{kernel/proc.c:/^kkill/}
允许一个进程请求终止另一个进程。
让
\lstinline{kill}
直接销毁目标进程会过于复杂，因为目标进程可能正在另一个 CPU 上执行，
也许正处于对内核数据结构进行敏感更新的中间状态。
因此
\lstinline{kkill}
做得很少：它只是设置目标进程的
\indexcode{p->killed}
标志，并且如果它在睡眠中，则将其唤醒。
最终，目标进程将进入或离开内核，
此时 \lstinline{usertrap} 中的代码会检查
\lstinline{p->killed}
是否被设置（通过调用
\lstinline{killed}
\lineref{kernel/proc.c:/^killed/} 来检查），
如果被设置，则调用 \lstinline{kexit}。
如果目标进程在用户空间运行，它将在下次通过系统调用或定时器（或其他设备）中断进入内核时发现自己已被杀死。

如果目标进程正在
\lstinline{sleep}
中，
\lstinline{kkill}
对
\lstinline{wakeup}
的调用将导致目标进程从
\lstinline{sleep}
返回。
这可能很危险，因为正在等待的条件可能不成立。
然而，xv6 中对
\lstinline{sleep}
的调用总是被包装在一个
\lstinline{while}
循环中，该循环在
\lstinline{sleep}
返回后重新测试条件。
有些对
\lstinline{sleep}
的调用还会在循环中测试
\lstinline{p->killed}，
如果被设置，则放弃当前活动。
这仅在放弃是正确的做法时才会执行。
例如，管道读写代码
\lineref{kernel/pipe.c:/killed.pr/}
在杀死标志被设置时返回；最终代码将返回到 trap，trap 会再次检查 \lstinline{p->killed} 并退出。

一些 xv6 的
\lstinline{sleep}
循环不会检查
\lstinline{p->killed}，
因为代码处于多步系统调用的中间，该调用应该是原子的
（即，如果在中途放弃会出错）。
virtio 驱动程序
\lineref{kernel/virtio\_disk.c:/sleep.b/}
就是一个例子：它不检查
\lstinline{p->killed}，
因为磁盘操作可能是一组写入操作中的一个，所有这些操作都需要完成，文件系统才能保持正确状态。
一个在等待磁盘 I/O 时被杀死的进程，直到完成当前系统调用并且
\lstinline{usertrap} 看到杀死标志时才会退出。

\section{进程锁}

与每个进程关联的锁（\lstinline{p->lock}）是 xv6 中最复杂的锁。
将 \lstinline{p->lock} 视为在读取或写入以下任何 \lstinline{struct proc} 字段时都必须持有的锁，是一种简单的思考方式：
\lstinline{p->state}、
\lstinline{p->chan}、
\lstinline{p->killed}、
\lstinline{p->xstate} 和
\lstinline{p->pid}。
这些字段可能被其他进程或其他 CPU 上的调度器线程使用，因此它们自然必须受到锁的保护。

然而，\lstinline{p->lock} 的大多数使用都是为了保护 xv6 进程数据结构和算法的更高级别不变量。以下是 \lstinline{p->lock} 的完整功能列表：

% 很难知道如何以统一的方式表述这些内容。
% 讨论应该围绕不变量？保护数据？避免竞争？
% 确保原子性？避免特定风险？

% “原子”的表述通常并不能真正说明风险所在。

\begin{itemize}

\item 与 \lstinline{p->state} 一起，防止在分配新进程的 \lstinline{proc[]} 槽位时发生竞争。
% 使检查 UNUSED 与分配原子化（或类似）。

\item 在进程被创建或销毁期间将其隐藏起来不被看到。
% 使分配和销毁的所有步骤原子化。

\item 防止父进程的 \lstinline{wait} 回收一个已将其状态设置为 \lstinline{ZOMBIE} 但尚未让出 CPU 的进程。
% 将状态设置为 ZOMBIE 和让出 CPU 原子化。

\item 防止另一个 CPU 的调度器在某个进程将状态设置为 \lstinline{RUNNABLE} 之后、完成 \lstinline{swtch} 之前，决定运行这个即将让出的进程。
% 将设置 p->state 和 swtch() 原子化。

\item 确保只有一个 CPU 的调度器决定运行一个 \lstinline{RUNNABLE} 的进程。
% 将调度器检查 RUNNABLE 与实际运行进程原子化。或者将状态设置为 RUNNING。

\item 防止定时器中断在进程处于 \lstinline{swtch} 时导致其让出 CPU。
% 使 swtch（实际上是整个序列）相对于定时器中断是原子的。

\item 与条件锁一起，防止 \lstinline{wakeup} 忽略一个正在调用 \lstinline{sleep} 但尚未完成让出 CPU 的进程。
% 避免 conditionCheck+sleep 与 wakeup 之间的竞争？
% 使条件检查与让出原子化？

\item 防止 \lstinline{kill} 的目标进程在 \lstinline{kkill} 检查 \lstinline{p->pid} 和设置 \lstinline{p->killed} 之间退出并可能被重新分配。
% 使 pid 检查与设置 killed 原子化。

\item 使 \lstinline{kkill} 对 \lstinline{p->state} 的检查和写入原子化。

\end{itemize}

\lstinline{p->parent} 字段受全局锁
\lstinline{wait_lock} 的保护，而不是 \lstinline{p->lock}。
只有进程的父进程会修改 \lstinline{p->parent}，
尽管该字段会被进程自身以及搜索其子进程的其他进程读取。
\lstinline{wait_lock} 的目的是当 \lstinline{wait} 睡眠等待任何子进程退出时充当条件锁。
退出的子进程在将其状态设置为 \lstinline{ZOMBIE}、唤醒其父进程并让出 CPU 之前，
会一直持有 \lstinline{wait_lock} 或 \lstinline{p->lock}。
\lstinline{wait_lock} 还序列化父进程和子进程并发的 \lstinline{exit}，
从而保证 \lstinline{init} 进程（它继承子进程）从其 \lstinline{wait} 中被唤醒。
\lstinline{wait_lock} 是一个全局锁，而不是每个父进程的每进程锁，
因为一个进程在获取它之前，无法知道其父进程是谁。

%% 
\section{现实世界}
%% 

\lstinline{sleep}
和
\lstinline{wakeup}
是一种简单而有效的同步方法，
但还有许多其他的方法；
信号量~\cite{dijkstra65} 就是一个例子。
所有方法中的第一个挑战是避免
我们在本章开头看到的“丢失唤醒”问题。
最初的 Unix 内核的
\lstinline{sleep}
只是简单地禁用中断，
这在 Unix 运行在单 CPU 系统上时是足够的。
由于 xv6 运行在多处理器上，
它为
\lstinline{sleep}
添加了一个显式的锁。
FreeBSD 的
\lstinline{msleep}
采用了相同的方法。
Plan 9 的
\lstinline{sleep}
使用了一个回调函数，该函数在即将进入睡眠之前持有调度锁的情况下运行；
该函数作为对睡眠条件的最后一刻检查，以避免丢失唤醒。
Linux 内核的
\lstinline{sleep}
使用一个显式的进程队列，称为等待队列，而不是等待通道；该队列有自己的内部锁。

在 \lstinline{wakeup} 中扫描整个进程表是低效的。
一个更好的解决方案是用一个数据结构替换
\lstinline{sleep}
和
\lstinline{wakeup}
中的
\lstinline{chan}，
该数据结构保存一个在该结构上睡眠的进程列表，
例如 Linux 的等待队列。
Plan 9 的
\lstinline{sleep}
和
\lstinline{wakeup}
称这种结构为会合点。
许多线程库将相同的结构称为条件变量；
在这种情况下，操作
\lstinline{sleep}
和
\lstinline{wakeup}
被称为
\lstinline{wait}
和
\lstinline{signal}。
所有这些机制都具有相同的风格：睡眠条件受某种锁的保护，该锁在睡眠期间原子地释放。

xv6 的
\lstinline{wakeup}
会唤醒在特定等待通道上等待的所有进程。
如果其中有多个进程，它们都会尝试获取条件锁并重新检查条件；
在许多情况下，只有一个进程能够执行有用的操作（例如，读取管道中等待的所有数据）。
其余的会发现条件不再成立，并重新进入睡眠状态；
唤醒它们浪费了 CPU 时间。
因此，大多数条件变量设计提供了两个原语：
\lstinline{signal}，
它唤醒等待条件变量的进程之一，以及
\lstinline{broadcast}，
它会唤醒所有等待的进程。

强制杀死进程会带来一些问题。
例如，一个被杀死的进程可能深度处于内核睡眠中，解开其堆栈需要小心，
因为调用堆栈上的每个函数可能都需要进行一些清理。
某些语言通过提供异常机制来提供帮助，但 C 语言不行。
此外，还有其他事件可能导致正在睡眠的进程被唤醒，即使它正在等待的事件尚未发生。
例如，当一个 Unix 进程处于睡眠状态时，另一个进程可能会向其发送一个
\indexcode{signal}。
在这种情况下，进程将从被中断的系统调用返回，返回值为 -1，错误码设置为 EINTR。
应用程序可以检查这些值并决定如何处理。
Xv6 不支持信号，因此不会出现这种复杂性。

Xv6 对 \lstinline{kill} 的支持并不完全令人满意：
有些睡眠循环可能应该检查 \lstinline{p->killed}。
一个相关的问题是，即使对于检查 \lstinline{p->killed} 的睡眠循环，
在 \lstinline{sleep} 和 \lstinline{kill} 之间也存在竞争；
后者可能刚好在受害者的循环检查完 \lstinline{p->killed} 之后、但在其调用 \lstinline{sleep} 之前设置 \lstinline{p->killed} 并尝试唤醒受害者。
如果发生此问题，受害者只有在它等待的条件发生时才会注意到 \lstinline{p->killed}。
这可能会延迟很久甚至永远不会发生（例如，如果受害者在等待控制台输入，而用户没有输入任何字符）。

%% 
\section{练习}
%% 

\begin{enumerate}

\item 在 xv6 中实现计数信号量。
选择 xv6 中几个使用 sleep 和 wakeup 的地方，用信号量替换它们。
评判结果。

\item 你能实现一个 {\tt sleep()} 的变体，它只接受一个参数（通道），而不需要锁参数吗？

\item 修复上述 \lstinline{kill} 和 \lstinline{sleep} 之间的竞争条件，
使得在受害者的睡眠循环检查完 \lstinline{p->killed} 之后、但在其调用 \lstinline{sleep} 之前发生的 \lstinline{kill} 能够导致受害者放弃当前系统调用。
% 答案：一种解决方案是在 sleep 中，在将进程状态设置为 sleep 之前检查 p->killed 是否被设置。

\item 设计一个方案，使每个睡眠循环都检查 \lstinline{p->killed}，
这样，例如，virtio 驱动程序中的进程如果被另一个进程杀死，就可以快速从 while 循环中返回。
% 答案：这很困难。现代 Unix 系统使用 setjmp 和 longjmp 以及非常仔细的编程来清理被中断的系统调用可能已构建的任何部分状态。

\end{enumerate}

  \caption{重叠锁以避免丢失唤醒}
    \label{fig:overlap}
\end{figure}

为什么 \lstinline{sleep} 和 \lstinline{wakeup} 的锁规则能确保即将进入睡眠的进程不会错过并发的唤醒？
即将进入睡眠的进程从{\it 在}检查条件之前，到{\it 在}将其自身标记为 \texttt{SLEEPING} 之后，
始终持有条件锁或其自身的 \lstinline{p->lock} 或两者都持有；
见图~\ref{fig:overlap}。
调用 \texttt{wakeup} 的进程需要获取{\it 两个}锁。
唤醒者可能先获取锁，这意味着它会在消费者线程检查条件之前使条件成立，
这样消费者线程就不需要调用 {\tt sleep()}；
或者唤醒者的 {\tt acquire()} 可能需要等待，直到消费者线程完全完成进入睡眠并释放锁，
此时唤醒者将看到消费者线程已被标记为 {\tt SLEEPING} 并将其唤醒。

有时多个进程在同一个通道上睡眠；例如，多个进程从一个管道读取数据。
一次对 \lstinline{wakeup} 的调用将唤醒它们全部。
其中一个将首先运行，并获取 \lstinline{sleep} 调用时所持有的锁，
并且（在管道的例子中）读取所有等待的数据。
其他进程会发现，尽管被唤醒了，但没有数据可读。
从它们的角度来看，这次唤醒是“虚假的”，它们必须再次进入睡眠。
因此，\lstinline{sleep} 总是在一个重新检查条件的循环内被调用，如上面的 \lstinline{P} 所示。

如果两个使用 sleep/wakeup 的地方意外选择了相同的通道，也不会造成伤害：
它们会遇到虚假唤醒，但如上所述的循环可以容忍这个问题。
sleep/wakeup 的魅力很大程度上在于它既轻量级（无需创建特殊的数据结构作为等待通道），
又提供了一层间接性（调用者无需知道他们正在与哪个特定进程交互）。

%% 
\section{代码：管道}
%% 
xv6 的管道是使用 \lstinline{sleep} 和 \lstinline{wakeup} 来同步生产者和消费者的代码示例。
我们在第~\ref{CH:UNIX}~章中了解了管道的接口：
写入管道一端的数据被复制到内核缓冲区，然后可以从另一端读取。
后续章节将研究围绕管道的文件描述符支持，现在让我们看看
\indexcode{pipewrite}
和
\indexcode{piperead}
的实现。

每个管道由一个
\indexcode{struct pipe}
表示，其中包含一个
\lstinline{lock}
和一个
\lstinline{data}
缓冲区。
字段
\lstinline{nread}
和
\lstinline{nwrite}
统计从缓冲区读取和写入缓冲区的总字节数。
缓冲区是环绕的：
在
\lstinline{buf[PIPESIZE-1]}
之后写入的下一个字节是
\lstinline{buf[0]}。
计数不会环绕。
这个约定使得实现能够区分缓冲区满
（\lstinline{nwrite}
\lstinline{==}
\lstinline{nread+PIPESIZE}）
和缓冲区空
（\lstinline{nwrite}
\lstinline{==}
\lstinline{nread}），
但这意味着对缓冲区的索引必须使用
\lstinline{buf[nread}
\lstinline{%}
\lstinline{PIPESIZE]}
而不是简单的
\lstinline{buf[nread]} 
（对 \lstinline{nwrite} 也同样处理）。

假设对
\lstinline{piperead}
和
\lstinline{pipewrite}
的调用在两个不同的 CPU 上同时发生。
\lstinline{pipewrite}
\lineref{kernel/pipe.c:/^pipewrite/}
首先获取管道的锁，该锁保护计数、数据及其相关不变量。
\lstinline{piperead}
\lineref{kernel/pipe.c:/^piperead/}
随后也尝试获取锁，但无法获取。
它会在
\lstinline{acquire}
\lineref{kernel/spinlock.c:/^acquire/}
中自旋，等待锁。
当
\lstinline{piperead}
等待时，
\lstinline{pipewrite}
循环处理正在写入的字节
（\lstinline{addr[0..n-1]}），
依次将每个字节添加到管道中
\lineref{kernel/pipe.c:/nwrite\+\+/}。
在这个循环过程中，可能会发生缓冲区填满的情况
\lineref{kernel/pipe.c:/DOC: pipewrite-full/}。
在这种情况下，
\lstinline{pipewrite}
调用
\lstinline{wakeup}
来提醒任何正在睡眠的读者缓冲区中有数据等待，
然后在
\lstinline{&pi->nwrite}
上睡眠，等待读者从缓冲区中取出一些字节。
\lstinline{sleep}
在将
\lstinline{pipewrite}
的进程置于睡眠状态时，会释放管道锁。

\lstinline{piperead}
现在获取了管道锁并进入其临界区：
它发现
\lstinline{pi->nread}
\lstinline{!=}
\lstinline{pi->nwrite}
\lineref{kernel/pipe.c:/DOC: pipe-empty/}
（\lstinline{pipewrite}
之所以进入睡眠，是因为
\lstinline{pi->nwrite}
\lstinline{==}
\lstinline{pi->nread}
\lstinline{+}
\lstinline{PIPESIZE}
\lineref{kernel/pipe.c:/pipewrite-full/}），
因此它进入
\lstinline{for}
循环，从管道中复制数据
\lineref{kernel/pipe.c:/DOC: piperead-copy/}，
并将
\lstinline{nread}
增加复制的字节数。
现在缓冲区中有了那么多可写入的空间，因此
\lstinline{piperead}
在返回之前调用
\lstinline{wakeup}
\lineref{kernel/pipe.c:/DOC: piperead-wakeup/}
来唤醒任何正在睡眠的写入者。
\lstinline{wakeup}
找到一个在
\lstinline{&pi->nwrite}
上睡眠的进程，即正在运行
\lstinline{pipewrite}
但因缓冲区满而停止的进程。
它将该进程标记为
\indexcode{RUNNABLE}。

管道代码为读者和写者使用了不同的等待通道
（\lstinline{pi->nread}
和
\lstinline{pi->nwrite}）；
在极少数情况下，如果有大量读者和写者等待同一个管道，这可能会提高系统效率。
管道代码在循环内睡眠，该循环会检查睡眠条件；
如果有多个读者或写者，除了第一个被唤醒的进程外，其他所有进程都会发现条件不成立，并再次进入睡眠。

%% 
\section{代码：Wait、Exit 和 Kill}

请阅读 {\tt kernel/proc.c} 中函数 {\tt kwait()}、{\tt kexit()} 和 {\tt kkill()} 的代码；
它们是相应系统调用的内部实现。

%% 
\lstinline{sleep}
和
\lstinline{wakeup}
可用于多种等待场景。
一个有趣的例子，在第~\ref{CH:UNIX}~章中介绍过，是子进程的 \indexcode{exit}
与其父进程的 \indexcode{wait} 之间的交互。
在子进程死亡时，父进程可能已经在 {\tt wait} 中睡眠，或者正在做其他事情；
在后一种情况下，后续对 {\tt wait} 的调用必须观察到子进程的死亡，
这可能发生在子进程调用 {\tt exit} 很久之后。
xv6 记录子进程的死亡直到 {\tt wait} 观察到它的方式是：
{\tt exit} 将调用者置于 \indexcode{ZOMBIE} 状态，
并保持该状态直到父进程的 {\tt wait} 注意到它，
然后将子进程的状态更改为 {\tt UNUSED}，复制子进程的退出状态，
并将子进程的进程 ID 返回给父进程。
如果父进程在子进程之前退出，父进程会将子进程交给
\lstinline{init}
进程，该进程会一直调用 {\tt wait}；
因此每个子进程都有一个父进程来为其进行清理。
一个挑战是避免同时发生的父进程
\lstinline{wait}
和子进程
\lstinline{exit}，
以及同时发生的
\lstinline{exit} 和 \lstinline{exit} 之间的竞争和死锁。

\lstinline{kwait}，即 {\tt wait} 的内核实现，
首先获取
\lstinline{wait_lock}
\lineref{kernel/proc.c:/^kwait/}，
它充当条件锁，有助于确保 \lstinline{kwait} 不会错过来自退出子进程的 \lstinline{wakeup}。
然后 \lstinline{kwait} 扫描进程表。
如果它找到一个处于 \texttt{ZOMBIE} 状态的子进程，
它会释放该子进程的资源及其
\lstinline{proc} 结构，
将子进程的退出状态复制到提供给 \lstinline{wait} 的地址（如果非 0），
并返回子进程的进程 ID。
如果
\lstinline{kwait}
找到了子进程，但都没有退出，
它会调用
\lstinline{sleep}
来等待其中任何一个退出
\lineref{kernel/proc.c:/DOC: wait-sleep/}，
然后再次扫描。
\lstinline{kwait} 通常持有两个锁，
\lstinline{wait_lock} 和某个进程的 \lstinline{pp->lock}；
避免死锁的顺序是先 \lstinline{wait_lock}，然后是 \lstinline{pp->lock}。

\lstinline{kexit} \lineref{kernel/proc.c:/^kexit/} 记录退出
状态，释放一些资源，调用 \lstinline{reparent} 将所有子进程交给 \lstinline{init} 进程，
唤醒可能正在 \lstinline{wait} 中的父进程，将调用者标记为僵尸进程，
并永久让出 CPU。
\lstinline{kexit} 在此序列期间同时持有
\lstinline{wait_lock} 和 \lstinline{p->lock}。
它持有 \lstinline{wait_lock} 是因为
它是 \lstinline{wakeup(p->parent)} 的条件锁，
防止父进程在 \lstinline{wait} 中丢失唤醒。
\lstinline{kexit} 在此序列中也必须持有
\lstinline{p->lock}，
以防止在 \lstinline{wait} 中的父进程在子进程最终调用 \lstinline{swtch} 之前就看到子进程处于
\lstinline{ZOMBIE} 状态。
\lstinline{kexit} 以与 \lstinline{kwait} 相同的顺序获取这些锁以避免死锁。

\lstinline{kexit} 在将其状态设置为 \lstinline{ZOMBIE} 之前唤醒父进程，这看起来可能不正确，
但这是安全的：
尽管
\lstinline{wakeup}
可能导致父进程运行，
但父进程的
\lstinline{kwait}
循环中的检查在子进程的 \lstinline{p->lock}
被 {\tt scheduler} 释放之前无法检查子进程的状态，
因此
\lstinline{kwait}
在
\lstinline{kexit}
将其状态设置为
\lstinline{ZOMBIE}
\lineref{kernel/proc.c:/state.=.ZOMBIE/}
之前无法查看退出的进程。

虽然
\lstinline{exit}
允许进程自行终止，
但 \lstinline{kill} 系统调用
\lineref{kernel/proc.c:/^kkill/}
允许一个进程请求终止另一个进程。
让
\lstinline{kill}
直接销毁目标进程会过于复杂，因为目标进程可能正在另一个 CPU 上执行，
也许正处于对内核数据结构进行敏感更新的中间状态。
因此
\lstinline{kkill}
做得很少：它只是设置目标进程的
\indexcode{p->killed}
标志，并且如果它在睡眠中，则将其唤醒。
最终，目标进程将进入或离开内核，
此时 \lstinline{usertrap} 中的代码会检查
\lstinline{p->killed}
是否被设置（通过调用
\lstinline{killed}
\lineref{kernel/proc.c:/^killed/} 来检查），
如果被设置，则调用 \lstinline{kexit}。
如果目标进程在用户空间运行，它将在下次通过系统调用或定时器（或其他设备）中断进入内核时发现自己已被杀死。

如果目标进程正在
\lstinline{sleep}
中，
\lstinline{kkill}
对
\lstinline{wakeup}
的调用将导致目标进程从
\lstinline{sleep}
返回。
这可能很危险，因为正在等待的条件可能不成立。
然而，xv6 中对
\lstinline{sleep}
的调用总是被包装在一个
\lstinline{while}
循环中，该循环在
\lstinline{sleep}
返回后重新测试条件。
有些对
\lstinline{sleep}
的调用还会在循环中测试
\lstinline{p->killed}，
如果被设置，则放弃当前活动。
这仅在放弃是正确的做法时才会执行。
例如，管道读写代码
\lineref{kernel/pipe.c:/killed.pr/}
在杀死标志被设置时返回；最终代码将返回到 trap，trap 会再次检查 \lstinline{p->killed} 并退出。

一些 xv6 的
\lstinline{sleep}
循环不会检查
\lstinline{p->killed}，
因为代码处于多步系统调用的中间，该调用应该是原子的
（即，如果在中途放弃会出错）。
virtio 驱动程序
\lineref{kernel/virtio\_disk.c:/sleep.b/}
就是一个例子：它不检查
\lstinline{p->killed}，
因为磁盘操作可能是一组写入操作中的一个，所有这些操作都需要完成，文件系统才能保持正确状态。
一个在等待磁盘 I/O 时被杀死的进程，直到完成当前系统调用并且
\lstinline{usertrap} 看到杀死标志时才会退出。

\section{进程锁}

与每个进程关联的锁（\lstinline{p->lock}）是 xv6 中最复杂的锁。
将 \lstinline{p->lock} 视为在读取或写入以下任何 \lstinline{struct proc} 字段时都必须持有的锁，是一种简单的思考方式：
\lstinline{p->state}、
\lstinline{p->chan}、
\lstinline{p->killed}、
\lstinline{p->xstate} 和
\lstinline{p->pid}。
这些字段可能被其他进程或其他 CPU 上的调度器线程使用，因此它们自然必须受到锁的保护。

然而，\lstinline{p->lock} 的大多数使用都是为了保护 xv6 进程数据结构和算法的更高级别不变量。以下是 \lstinline{p->lock} 的完整功能列表：

% 很难知道如何以统一的方式表述这些内容。
% 讨论应该围绕不变量？保护数据？避免竞争？
% 确保原子性？避免特定风险？

% “原子”的表述通常并不能真正说明风险所在。

\begin{itemize}

\item 与 \lstinline{p->state} 一起，防止在分配新进程的 \lstinline{proc[]} 槽位时发生竞争。
% 使检查 UNUSED 与分配原子化（或类似）。

\item 在进程被创建或销毁期间将其隐藏起来不被看到。
% 使分配和销毁的所有步骤原子化。

\item 防止父进程的 \lstinline{wait} 回收一个已将其状态设置为 \lstinline{ZOMBIE} 但尚未让出 CPU 的进程。
% 将状态设置为 ZOMBIE 和让出 CPU 原子化。

\item 防止另一个 CPU 的调度器在某个进程将状态设置为 \lstinline{RUNNABLE} 之后、完成 \lstinline{swtch} 之前，决定运行这个即将让出的进程。
% 将设置 p->state 和 swtch() 原子化。

\item 确保只有一个 CPU 的调度器决定运行一个 \lstinline{RUNNABLE} 的进程。
% 将调度器检查 RUNNABLE 与实际运行进程原子化。或者将状态设置为 RUNNING。

\item 防止定时器中断在进程处于 \lstinline{swtch} 时导致其让出 CPU。
% 使 swtch（实际上是整个序列）相对于定时器中断是原子的。

\item 与条件锁一起，防止 \lstinline{wakeup} 忽略一个正在调用 \lstinline{sleep} 但尚未完成让出 CPU 的进程。
% 避免 conditionCheck+sleep 与 wakeup 之间的竞争？
% 使条件检查与让出原子化？

\item 防止 \lstinline{kill} 的目标进程在 \lstinline{kkill} 检查 \lstinline{p->pid} 和设置 \lstinline{p->killed} 之间退出并可能被重新分配。
% 使 pid 检查与设置 killed 原子化。

\item 使 \lstinline{kkill} 对 \lstinline{p->state} 的检查和写入原子化。

\end{itemize}

\lstinline{p->parent} 字段受全局锁
\lstinline{wait_lock} 的保护，而不是 \lstinline{p->lock}。
只有进程的父进程会修改 \lstinline{p->parent}，
尽管该字段会被进程自身以及搜索其子进程的其他进程读取。
\lstinline{wait_lock} 的目的是当 \lstinline{wait} 睡眠等待任何子进程退出时充当条件锁。
退出的子进程在将其状态设置为 \lstinline{ZOMBIE}、唤醒其父进程并让出 CPU 之前，
会一直持有 \lstinline{wait_lock} 或 \lstinline{p->lock}。
\lstinline{wait_lock} 还序列化父进程和子进程并发的 \lstinline{exit}，
从而保证 \lstinline{init} 进程（它继承子进程）从其 \lstinline{wait} 中被唤醒。
\lstinline{wait_lock} 是一个全局锁，而不是每个父进程的每进程锁，
因为一个进程在获取它之前，无法知道其父进程是谁。

%% 
\section{现实世界}
%% 

\lstinline{sleep}
和
\lstinline{wakeup}
是一种简单而有效的同步方法，
但还有许多其他的方法；
信号量~\cite{dijkstra65} 就是一个例子。
所有方法中的第一个挑战是避免
我们在本章开头看到的“丢失唤醒”问题。
最初的 Unix 内核的
\lstinline{sleep}
只是简单地禁用中断，
这在 Unix 运行在单 CPU 系统上时是足够的。
由于 xv6 运行在多处理器上，
它为
\lstinline{sleep}
添加了一个显式的锁。
FreeBSD 的
\lstinline{msleep}
采用了相同的方法。
Plan 9 的
\lstinline{sleep}
使用了一个回调函数，该函数在即将进入睡眠之前持有调度锁的情况下运行；
该函数作为对睡眠条件的最后一刻检查，以避免丢失唤醒。
Linux 内核的
\lstinline{sleep}
使用一个显式的进程队列，称为等待队列，而不是等待通道；该队列有自己的内部锁。

在 \lstinline{wakeup} 中扫描整个进程表是低效的。
一个更好的解决方案是用一个数据结构替换
\lstinline{sleep}
和
\lstinline{wakeup}
中的
\lstinline{chan}，
该数据结构保存一个在该结构上睡眠的进程列表，
例如 Linux 的等待队列。
Plan 9 的
\lstinline{sleep}
和
\lstinline{wakeup}
称这种结构为会合点。
许多线程库将相同的结构称为条件变量；
在这种情况下，操作
\lstinline{sleep}
和
\lstinline{wakeup}
被称为
\lstinline{wait}
和
\lstinline{signal}。
所有这些机制都具有相同的风格：睡眠条件受某种锁的保护，该锁在睡眠期间原子地释放。

xv6 的
\lstinline{wakeup}
会唤醒在特定等待通道上等待的所有进程。
如果其中有多个进程，它们都会尝试获取条件锁并重新检查条件；
在许多情况下，只有一个进程能够执行有用的操作（例如，读取管道中等待的所有数据）。
其余的会发现条件不再成立，并重新进入睡眠状态；
唤醒它们浪费了 CPU 时间。
因此，大多数条件变量设计提供了两个原语：
\lstinline{signal}，
它唤醒等待条件变量的进程之一，以及
\lstinline{broadcast}，
它会唤醒所有等待的进程。

强制杀死进程会带来一些问题。
例如，一个被杀死的进程可能深度处于内核睡眠中，解开其堆栈需要小心，
因为调用堆栈上的每个函数可能都需要进行一些清理。
某些语言通过提供异常机制来提供帮助，但 C 语言不行。
此外，还有其他事件可能导致正在睡眠的进程被唤醒，即使它正在等待的事件尚未发生。
例如，当一个 Unix 进程处于睡眠状态时，另一个进程可能会向其发送一个
\indexcode{signal}。
在这种情况下，进程将从被中断的系统调用返回，返回值为 -1，错误码设置为 EINTR。
应用程序可以检查这些值并决定如何处理。
Xv6 不支持信号，因此不会出现这种复杂性。

Xv6 对 \lstinline{kill} 的支持并不完全令人满意：
有些睡眠循环可能应该检查 \lstinline{p->killed}。
一个相关的问题是，即使对于检查 \lstinline{p->killed} 的睡眠循环，
在 \lstinline{sleep} 和 \lstinline{kill} 之间也存在竞争；
后者可能刚好在受害者的循环检查完 \lstinline{p->killed} 之后、但在其调用 \lstinline{sleep} 之前设置 \lstinline{p->killed} 并尝试唤醒受害者。
如果发生此问题，受害者只有在它等待的条件发生时才会注意到 \lstinline{p->killed}。
这可能会延迟很久甚至永远不会发生（例如，如果受害者在等待控制台输入，而用户没有输入任何字符）。

%% 
\section{练习}
%% 

\begin{enumerate}

\item 在 xv6 中实现计数信号量。
选择 xv6 中几个使用 sleep 和 wakeup 的地方，用信号量替换它们。
评判结果。

\item 你能实现一个 {\tt sleep()} 的变体，它只接受一个参数（通道），而不需要锁参数吗？

\item 修复上述 \lstinline{kill} 和 \lstinline{sleep} 之间的竞争条件，
使得在受害者的睡眠循环检查完 \lstinline{p->killed} 之后、但在其调用 \lstinline{sleep} 之前发生的 \lstinline{kill} 能够导致受害者放弃当前系统调用。
% 答案：一种解决方案是在 sleep 中，在将进程状态设置为 sleep 之前检查 p->killed 是否被设置。

\item 设计一个方案，使每个睡眠循环都检查 \lstinline{p->killed}，
这样，例如，virtio 驱动程序中的进程如果被另一个进程杀死，就可以快速从 while 循环中返回。
% 答案：这很困难。现代 Unix 系统使用 setjmp 和 longjmp 以及非常仔细的编程来清理被中断的系统调用可能已构建的任何部分状态。

\end{enumerate}
